% Luc Destombe --- licence CC BY-NC-SA 4.0
% https://creativecommons.org/licenses/by-nc-sa/4.0/deed.fr
% Fichier autonome : compiler avec pdflatex, deux fois (signets, nombre de pages).
\documentclass[11pt,a4paper]{article}
\makeatletter
\newif\iflycee@unecolonne
\lycee@unecolonnetrue
\newif\iflycee@palatino
\lycee@palatinotrue

\RequirePackage[utf8]{inputenc}
\RequirePackage[T1]{fontenc}
\RequirePackage[french]{babel}
\RequirePackage{etoolbox}

\newcommand{\ShorthandsOff}{\@ifundefined{shorthandoff}{}{\shorthandoff{;:!?}}}
\newcommand{\ShorthandsOn}{\@ifundefined{shorthandon}{}{\shorthandon{;:!?}}}
\ShorthandsOff
\AtBeginDocument{\ShorthandsOn}
\AtBeginEnvironment{tikzpicture}{\ShorthandsOff}

\RequirePackage{geometry}
\RequirePackage{fancyhdr}
\RequirePackage{lastpage}
\RequirePackage{multicol}
\RequirePackage{array}
\RequirePackage{enumitem}
\RequirePackage{xcolor}
\RequirePackage{needspace}

\RequirePackage{amsmath,amssymb}
\RequirePackage{mathpazo}
\DeclareMathSizes{10.95}{10.95}{8}{5.5}
\begingroup\catcode`\_=\active \catcode`\^=\active
\gdef_#1{\sb{\scriptscriptstyle #1}}
\gdef^#1{\sp{\scriptscriptstyle\lycee@moinsepais #1}}
\endgroup
\newcommand\lycee@moins{\mathbin{%
  \setbox\z@\hbox{$\m@th\scriptscriptstyle\lycee@moinsorig$}%
  \dimen@=.55\wd\z@ \advance\dimen@-.5pt
  \hbox to.75\wd\z@{\hss$\m@th\scriptscriptstyle\vcenter{\hbox{%
    \tikz\draw[line width=.5pt,line cap=round](0,0)--(\the\dimen@,0);}}$\hss}}}
\begingroup\lccode`\~=`\- \lowercase{\endgroup\let~}\lycee@moins
\newcommand\lycee@moinsepais{\mathcode`\-="8000 }
\AtBeginDocument{\mathchardef\lycee@moinsorig=\mathcode`\-\relax}
\AtBeginDocument{\catcode`\_=12 \mathcode`\_="8000
  \catcode`\^=12 \mathcode`\^="8000 }
\newcommand\lycee@exposants{%
  \fontdimen13\textfont2=.48\fontdimen6\textfont2
  \fontdimen14\textfont2=.48\fontdimen6\textfont2
  \fontdimen15\textfont2=.48\fontdimen6\textfont2
  \fontdimen13\scriptfont2=.48\fontdimen6\scriptfont2
  \fontdimen14\scriptfont2=.48\fontdimen6\scriptfont2
  \fontdimen15\scriptfont2=.48\fontdimen6\scriptfont2}
\every@math@size\expandafter{\the\every@math@size\lycee@exposants}
\newlength{\retraitcalculs}
\setlength{\retraitcalculs}{2em}

\RequirePackage{tikz}
\usetikzlibrary{arrows.meta,calc,babel,shapes.misc}
\RequirePackage[most]{tcolorbox}
\tcbuselibrary{skins,breakable}

\definecolor{coulDef}{HTML}{1F4E79}
\definecolor{coulProp}{HTML}{7030A0}
\definecolor{coulRem}{HTML}{C55A11}
\definecolor{coulEx}{HTML}{2E7D32}
\definecolor{coulExo}{HTML}{B03A2E}
\definecolor{coulPart}{HTML}{1F4E79}
\definecolor{coulMeth}{HTML}{1B6E4F}
\definecolor{coulCourbe}{HTML}{0B5ED7}
\definecolor{coulCourbeB}{HTML}{C00000}
\definecolor{coulCourbeC}{HTML}{2E7D32}
\definecolor{coulGrille}{HTML}{8C8C8C}
\definecolor{coulRep}{HTML}{1B6E4F}
\definecolor{coulBar}{HTML}{2E7D32}

\newcommand{\R}{\mathbb{R}}
\newcommand{\Rs}{\mathbb{R}^{*}}
\renewcommand{\le}{\leqslant}
\renewcommand{\ge}{\geqslant}
\renewcommand{\approx}{\simeq}
\newcommand{\vect}[1]{\overrightarrow{#1}}
\newcommand{\norme}[1]{\left\lVert #1 \right\rVert}
\newcommand{\intervalle}[2]{\left[#1\,;#2\right]}
\RequirePackage{cancel}
\renewcommand{\CancelColor}{\color{coulExo}}
\newcommand{\fois}[1]{\times{\color{coulExo}#1}}
\newcommand{\barre}[1]{\cancel{#1}}

\tikzset{grille/.style={coulGrille,line width=0.4pt}}
\newcommand{\quadrillage}[4]{\draw[grille,step=1] (#1,#2) grid (#3,#4);}
\tikzset{
  croix/.style={cross out,draw,line width=0.6pt,inner sep=0pt,minimum size=4pt},
  croixdroite/.style={croix,rotate=45,line width=0.8pt},
}
\newcommand{\pt}[3]{\path (#1) node[croix]{} node[#2,font=\small,inner sep=2.5pt]{$#3$};}

\newcommand{\lycee@repere}[4]{%
  \draw[grille,step=1] (#1,#2) grid (#3,#4);
  \draw[-{Stealth[length=2mm]},thick] (#1,0)--({#3+0.4},0) node[below left=-1pt]{$x$};
  \draw[-{Stealth[length=2mm]},thick] (0,#2)--(0,{#4+0.4}) node[below left=-1pt]{$y$};
  \node[below left,font=\scriptsize,inner sep=1.5pt] at (0,0) {$0$};}
\newcommand{\lycee@gradx}[1]{\foreach \k/\l in {#1}{%
  \draw (\k,0.07)--(\k,-0.07) node[below,font=\scriptsize,inner sep=1.5pt]{$\l$};}}
\newcommand{\lycee@grady}[1]{\foreach \k/\l in {#1}{%
  \draw (0.07,\k)--(-0.07,\k) node[left,font=\scriptsize,inner sep=1.5pt]{$\l$};}}
\newcommand{\lycee@intff}[2]{\left[#1\,;#2\right]}
\newcommand{\lycee@intfo}[2]{\left[#1\,;#2\right[}
\newcommand{\lycee@intof}[2]{\left]#1\,;#2\right]}
\newcommand{\lycee@intoo}[2]{\left]#1\,;#2\right[}
\newcommand{\lycee@ens}[1]{\left\{#1\right\}}
\AtBeginDocument{%
  \providecommand{\repere}{\lycee@repere}%
  \providecommand{\gradx}{\lycee@gradx}%
  \providecommand{\grady}{\lycee@grady}%
  \providecommand{\Cf}{\mathcal{C}\sb{\scriptscriptstyle f}}%
  \providecommand{\Cg}{\mathcal{C}\sb{\scriptscriptstyle g}}%
  \providecommand{\intff}{\lycee@intff}%
  \providecommand{\intfo}{\lycee@intfo}%
  \providecommand{\intof}{\lycee@intof}%
  \providecommand{\intoo}{\lycee@intoo}%
  \providecommand{\ens}{\lycee@ens}}

\newlist{questions}{enumerate}{3}
\setlist[questions,1]{label=\arabic*\textdegree),leftmargin=*,itemsep=2pt,topsep=3pt}
\setlist[questions,2]{label=\alph*),leftmargin=*,itemsep=1pt,topsep=2pt}
\setlist[questions,3]{label=\roman*),leftmargin=*,itemsep=1pt,topsep=1pt}
\setlist[itemize]{label=\textbullet,leftmargin=*,itemsep=2pt,topsep=3pt}

\newcommand{\besoinplace}[1]{\par\penalty\@M\begingroup
  \dimen@=#1\relax\dimen@ii=\pagegoal\advance\dimen@ii-\pagetotal
  \ifdim\dimen@>\dimen@ii\ifdim\dimen@ii>\z@\vfil\fi\break\fi\endgroup}

\newcommand{\lycee@pied}{\thepage/\pageref{LastPage}}
\newcommand{\entete}[2]{%
  \pagestyle{fancy}\fancyhf{}%
  \fancyhead[L]{\footnotesize\color{coulPart}\textbf{#1}}%
  \fancyhead[R]{\footnotesize\color{coulPart}\textbf{#2}}%
  \fancyfoot[C]{\footnotesize\lycee@pied}%
  \renewcommand{\headrulewidth}{0.4pt}%
  \renewcommand{\headrule}{\hbox to\headwidth{%
    \color{coulPart!40}\leaders\hrule height \headrulewidth\hfill}}}

\RequirePackage{ccicons}
\newcommand{\lycee@licence}{{\scriptsize\color{gray}%
  \copyright\ Luc Destombe --- \ccbyncsa\ CC BY-NC-SA 4.0}}
\newif\iflycee@licence \lycee@licencetrue
\newcommand{\sanslicence}{\lycee@licencefalse}
\AtBeginDocument{\iflycee@licence\AddToHookNext{shipout/foreground}{%
  \put(0,0){\rlap{\kern\dimexpr1in+\hoffset+\oddsidemargin+\textwidth\relax
    \llap{\raisebox{-\dimexpr1in+\voffset+\topmargin+\headheight+\headsep
      +\textheight+\footskip\relax}{\lycee@licence}}}}}\fi}

\newcommand{\formulesserrees}{%
  \setlength{\abovedisplayskip}{3pt}\setlength{\belowdisplayskip}{3pt}%
  \setlength{\abovedisplayshortskip}{2pt}\setlength{\belowdisplayshortskip}{3pt}}

\setlength{\parindent}{0pt}

\RequirePackage[implicit=false,hidelinks,pdfencoding=auto,psdextra,bookmarksopen,
  bookmarksopenlevel=1,pdfcreator={},pdfproducer={}]{hyperref}
\RequirePackage{bookmark}
\hypersetup{pdfauthor={Luc Destombe},pdfkeywords={CC BY-NC-SA 4.0}}
\newcounter{lycee@signet}
\newcommand{\signet}[3][]{\stepcounter{lycee@signet}%
  \edef\lycee@dest{\if\relax\detokenize{#1}\relax
    signet.\arabic{lycee@signet}\else#1\fi}%
  \hypertarget{\lycee@dest}{}\bookmark[level=#2,dest=\lycee@dest]{#3}}
\pdfstringdefDisableCommands{%
  \def\R{R}\def\vect#1{#1}\def\Cf{Cf}\def\Cg{Cg}\def\fois#1{×#1}%
  \def\barre#1{#1}\def\intervalle#1#2{[#1;#2]}\def\intff#1#2{[#1;#2]}%
  \def\textsuperscript#1{#1}\def\dfrac#1#2{#1/#2}\def\frac#1#2{#1/#2}%
  \def\vec#1{#1⃗}}%
\RequirePackage[official]{eurosym}

\tcbset{exercice corrige/.style={
  enhanced,breakable,before upper=\raggedright,
  colback=white,colframe=coulExo,
  boxrule=0.8pt,arc=2pt,
  left=6pt,right=6pt,top=8pt,bottom=5pt,
  before skip=9pt,after skip=4pt,
  coltitle=white,fonttitle=\bfseries\small,
  attach boxed title to top left={xshift=8pt,yshift=-\tcboxedtitleheight/2},
  boxed title style={colback=coulExo,arc=2pt,boxrule=0pt}}}
\newcounter{exo}
\newtcolorbox{exercice}[1][]{exercice corrige,
  phantom={\signet[exercice-\arabic{exo}]{2}{Exercice \arabic{exo}}},
  title={Exercice \theexo\ifx\relax#1\relax\else{} --- #1\fi}}
\newenvironment{exo}[1][\relax]{\refstepcounter{exo}\begin{exercice}[#1]}{\end{exercice}}

\geometry{top=1.8cm,bottom=1cm,left=1cm,right=1cm,headheight=14pt,footskip=0.6cm}
\renewcommand{\lycee@pied}{\thepage}

\renewcommand{\theexo}{\ifnum\value{exo}<10 0\fi\arabic{exo}}
\tcbset{exercice corrige/.style={
  enhanced,breakable,before upper=\raggedright,
  colback=white,colframe=coulExo,
  boxrule=0.9pt,arc=2pt,
  left=8pt,right=8pt,top=8pt,bottom=6pt,
  before skip=8pt,after skip=6pt,
  coltitle=white,fonttitle=\bfseries,
  attach boxed title to top left={xshift=8pt,yshift=-\tcboxedtitleheight/2},
  boxed title style={colback=coulExo,arc=2pt,boxrule=0pt}}}

\newcommand{\entetefiche}[2]{%
  \begin{center}
    \begin{tcolorbox}[enhanced,width=0.92\linewidth,colback=white,colframe=coulPart,
        boxrule=1.4pt,arc=3pt,halign=center,top=6pt,bottom=6pt]
      {\large\bfseries\color{coulPart} CORRIGÉ --- FICHE D'EXERCICES}\\[3pt]
      {\normalsize\bfseries #1}\\[2pt]
      {\normalsize #2}
    \end{tcolorbox}
  \end{center}}

\newcounter{partie}
\newcommand{\partie}[1]{%
  \stepcounter{partie}%
  \besoinplace{8\baselineskip}\vspace{4pt}%
  \begin{tcolorbox}[enhanced,colback=coulPart!8,colframe=coulPart,boxrule=0pt,
      leftrule=3pt,arc=0pt,left=6pt,right=4pt,top=3pt,bottom=3pt,
      before skip=6pt,after skip=2pt]
    \bfseries\color{coulPart}\Roman{partie} --- #1\signet{1}{\Roman{partie} --- #1}
  \end{tcolorbox}}
\newcommand{\partiesansnum}[1]{%
  \besoinplace{8\baselineskip}\vspace{4pt}%
  \begin{tcolorbox}[enhanced,colback=coulPart!8,colframe=coulPart,boxrule=0pt,
      leftrule=3pt,arc=0pt,left=6pt,right=4pt,top=3pt,bottom=3pt,
      before skip=6pt,after skip=2pt]
    \bfseries\color{coulPart}#1\signet{1}{#1}
  \end{tcolorbox}}

\setlist[questions,1]{label=\arabic*\textdegree),leftmargin=*,itemsep=3pt,topsep=3pt}
\setlist[questions,2]{label=\alph*),leftmargin=*,itemsep=2pt,topsep=2pt}
\setlist[questions,3]{label=\roman*),leftmargin=*,itemsep=1pt,topsep=1pt}
\newlist{expr}{enumerate}{1}
\setlist[expr]{label={\Alph*\ $=$},leftmargin=*,itemsep=2pt,topsep=3pt}

\newcommand{\res}[1]{{\color{coulRep}#1}}
\newcommand{\rem}[1]{\par\smallskip{\small\itshape\color{black!60}$\rightarrow$ #1}}
\newenvironment{calculs}{\par\smallskip\noindent$\displaystyle\begin{aligned}[t]}%
  {\end{aligned}$\par\smallskip}
\newcommand{\justif}[1]{\quad\text{\small\itshape\color{black!60}#1}}
\newcommand{\cas}[1]{\par\smallskip\noindent\textbf{\color{coulExo}#1}\nobreak}
\newcommand{\explique}[2]{\par\smallskip\noindent
  \begin{minipage}[t]{0.6\linewidth}\raggedright #1\end{minipage}\hfill
  \begin{minipage}[t]{0.37\linewidth}\raggedright\leavevmode
    {\small\itshape\color{black!60}#2}\end{minipage}\par}

\setlength{\parskip}{2pt}

\makeatother
\usepackage{tkz-tab}
\entete{Chapitre 3 --- Polynômes du second degré --- corrigé}{1\textsuperscript{re} STI2D}
\setlength{\parindent}{0pt}
\setlength{\jot}{4pt}

\let\casligne\cas
\renewcommand{\cas}[1]{\casligne{#1}\par\nopagebreak}

\newcommand{\tabsignes}[4][2.2]{\par\smallskip{\centering
  \begin{tikzpicture}[scale=0.85]
  \tkzTabInit[lgt=2.6,espcl=#1]{$x$/1,$#2$/1.2}{#3}
  \tkzTabLine{#4}
  \end{tikzpicture}\par}\smallskip}
\newcommand{\tabvar}[4][2.2]{\par\smallskip{\centering
  \begin{tikzpicture}[scale=0.85]
  \tkzTabInit[lgt=2,espcl=#1]{$x$/1,$#2$/2}{#3}
  \tkzTabVar{#4}
  \end{tikzpicture}\par}\smallskip}

\definecolor{coulInt}{HTML}{D62828}
\definecolor{coulIntB}{HTML}{1565C0}
\NewDocumentCommand{\morceau}{O{coulInt} m m}{%
  \draw[#1,line width=2.2pt,opacity=0.85] (#2,0)--(#3,0);}

\begin{document}

\entetefiche{Chapitre 3 --- Polynômes du second degré}{1\textsuperscript{re} STI2D}

\partie{Fonction polynôme du second degré}

\begin{exo}[Est-ce un trinôme ?]
\explique{Une fonction polynôme du second degré s'écrit $ax^{2}+bx+c$ avec $a\ne 0$.
Quand il y a des parenthèses, on développe d'abord.}{méthode 1 du cours}
\cas{a)}
\explique{$f(x)=3x^{2}-2x+5$ : \res{oui}, avec $a=\res{3}$, $b=\res{-2}$ et
$c=\res{5}$.}{chaque coefficient avec son signe}
\cas{b)}
\explique{$g(x)=4x+1$ : \res{non}, il n'y a pas de terme en $x^{2}$ ; c'est une fonction
affine.}{$a=0$}
\cas{c)}
\explique{$h(x)=-x^{2}+6$ : \res{oui}, avec $a=\res{-1}$, $b=\res{0}$ et
$c=\res{6}$.}{$-x^{2}=-1\times x^{2}$ ; pas de terme en $x$ : $b=0$}
\cas{d)}
\explique{$k(x)=2x^{3}+x^{2}-1$ : \res{non}, il y a un terme en $x^{3}$.}{le plus grand
exposant doit être $2$}
\cas{e)}
\begin{calculs}
l(x) &= (x+2)(x-4) &&\justif{des parenthèses : on développe}\\
l(x) &= x^{2}-4x+2x-8 &&\justif{quatre produits}\\
l(x) &= x^{2}-2x-8 &&\justif{on réduit : $-4x+2x=-2x$}
\end{calculs}
\explique{\res{Oui}, avec $a=\res{1}$, $b=\res{-2}$ et $c=\res{-8}$.}{}
\cas{f)}
\begin{calculs}
m(x) &= x^{2}-(x-1)^{2} &&\justif{on recopie}\\
m(x) &= x^{2}-\left(x^{2}-2x+1\right) &&\justif{identité remarquable, entre parenthèses}\\
m(x) &= x^{2}-x^{2}+2x-1 &&\justif{le $-$ devant la parenthèse change tous les signes}\\
m(x) &= 2x-1 &&\justif{les $x^{2}$ s'éliminent}
\end{calculs}
\explique{\res{Non} : $m$ est une fonction affine.}{$a=0$}
\rem{Erreur fréquente en f) : écrire $x^{2}-x^{2}-2x+1$ en oubliant de changer tous les
signes de la parenthèse.}
\end{exo}

\begin{exo}[Lire les coefficients]
\explique{On lit chaque coefficient \emph{avec son signe} ; un terme absent a un
coefficient nul.}{méthode 1 du cours}
\cas{$f(x)=3x^{2}-4x+2$}
\explique{$a=\res{3}$, $b=\res{-4}$, $c=\res{2}$.}{}
\cas{$g(x)=-x^{2}+5x$}
\explique{$a=\res{-1}$, $b=\res{5}$, $c=\res{0}$.}{$-x^{2}=-1\times x^{2}$ ; pas de
nombre seul : $c=0$}
\cas{$h(x)=x^{2}-16$}
\explique{$a=\res{1}$, $b=\res{0}$, $c=\res{-16}$.}{$x^{2}=1\times x^{2}$ ; pas de terme
en $x$ : $b=0$}
\cas{$k(x)=-2x^{2}+7$}
\explique{$a=\res{-2}$, $b=\res{0}$, $c=\res{7}$.}{}
\cas{$l(x)=\frac{1}{3}x^{2}+x-4$}
\explique{$a=\res{\dfrac{1}{3}}$, $b=\res{1}$, $c=\res{-4}$.}{$x=1\times x$}
\cas{$m(x)=6-3x-x^{2}$}
\explique{On range d'abord : $m(x)=-x^{2}-3x+6$. Donc $a=\res{-1}$, $b=\res{-3}$,
$c=\res{6}$.}{$a$ est devant $x^{2}$, quelle que soit sa place}
\cas{$n(x)=-x^{2}+2x-5$}
\explique{$a=\res{-1}$, $b=\res{2}$, $c=\res{-5}$.}{}
\cas{$p(x)=5x^{2}$}
\explique{$a=\res{5}$, $b=\res{0}$, $c=\res{0}$.}{c'est bien un trinôme : seul $a$ doit
être non nul}
\rem{Erreur fréquente pour $m$ : prendre $a=6$ parce que $6$ est écrit en premier.}
\end{exo}

\begin{exo}[Développer d'abord]
\explique{On développe, on réduit et on range : $x^{2}$, puis $x$, puis le nombre seul.
Si les $x^{2}$ s'éliminent, ce n'est pas un trinôme.}{méthode 1 du cours}
\cas{$f$}
\begin{calculs}
f(x) &= (x+3)(x-4) &&\justif{on recopie}\\
f(x) &= x^{2}-4x+3x-12 &&\justif{quatre produits}\\
f(x) &= \res{x^{2}-x-12} &&\justif{$-4x+3x=-x$}
\end{calculs}
\explique{Trinôme : $a=\res{1}$, $b=\res{-1}$, $c=\res{-12}$.}{}
\cas{$g$}
\begin{calculs}
g(x) &= (3x-2)(x+1) &&\justif{on recopie}\\
g(x) &= 3x^{2}+3x-2x-2 &&\justif{quatre produits}\\
g(x) &= \res{3x^{2}+x-2} &&\justif{$3x-2x=x$}
\end{calculs}
\explique{Trinôme : $a=\res{3}$, $b=\res{1}$, $c=\res{-2}$.}{}
\cas{$h$}
\begin{calculs}
h(x) &= (x-2)^{2}+5 &&\justif{on recopie}\\
h(x) &= x^{2}-4x+4+5 &&\justif{$(a-b)^{2}=a^{2}-2ab+b^{2}$}\\
h(x) &= \res{x^{2}-4x+9} &&
\end{calculs}
\explique{Trinôme : $a=\res{1}$, $b=\res{-4}$, $c=\res{9}$.}{}
\cas{$k$}
\begin{calculs}
k(x) &= 2(x+3)^{2}-1 &&\justif{on recopie}\\
k(x) &= 2\left(x^{2}+6x+9\right)-1 &&\justif{le carré d'abord, entre parenthèses}\\
k(x) &= 2x^{2}+12x+18-1 &&\justif{puis on multiplie par $2$}\\
k(x) &= \res{2x^{2}+12x+17} &&
\end{calculs}
\explique{Trinôme : $a=\res{2}$, $b=\res{12}$, $c=\res{17}$.}{}
\cas{$l$}
\begin{calculs}
l(x) &= (x+3)^{2}-(x-3)^{2} &&\justif{on recopie}\\
l(x) &= x^{2}+6x+9-\left(x^{2}-6x+9\right) &&\justif{deux identités remarquables}\\
l(x) &= x^{2}+6x+9-x^{2}+6x-9 &&\justif{le $-$ change tous les signes}\\
l(x) &= \res{12x} &&
\end{calculs}
\explique{\res{Pas un trinôme} : les $x^{2}$ s'éliminent.}{$a=0$}
\cas{$m$}
\begin{calculs}
m(x) &= (3x+1)^{2}-9x^{2} &&\justif{on recopie}\\
m(x) &= 9x^{2}+6x+1-9x^{2} &&\justif{$(3x)^{2}=9x^{2}$}\\
m(x) &= \res{6x+1} &&
\end{calculs}
\explique{\res{Pas un trinôme} : les $x^{2}$ s'éliminent.}{$a=0$}
\rem{Erreur fréquente : $(x-2)^{2}=x^{2}+4$. Le carré d'une différence a trois termes,
dont le double produit $-4x$.}
\end{exo}

\begin{exo}[Calculer des images]
\explique{On remplace $x$ par le nombre, entre parenthèses s'il est négatif, puis on
calcule : les puissances d'abord, puis les produits.}{priorités de calcul}
\cas{1)}
\begin{calculs}
f(0) &= 3\times 0^{2}-2\times 0+1 &&\justif{on remplace $x$ par $0$}\\
f(0) &= \res{1} &&\justif{l'image de $0$ est toujours $c$}\\[3pt]
f(1) &= 3\times 1^{2}-2\times 1+1 &&\\
f(1) &= 3-2+1=\res{2} &&\\[3pt]
f(-2) &= 3\times(-2)^{2}-2\times(-2)+1 &&\justif{$-2$ entre parenthèses}\\
f(-2) &= 12+4+1=\res{17} &&\justif{$(-2)^{2}=4$ et $-2\times(-2)=+4$}\\[3pt]
f\!\left(\tfrac{1}{3}\right) &= 3\times\left(\frac{1}{3}\right)^{2}-2\times\frac{1}{3}+1 &&\\
f\!\left(\tfrac{1}{3}\right) &= 3\times\frac{1}{9}-\frac{2}{3}+1 &&\justif{$\left(\frac{1}{3}\right)^{2}=\frac{1}{9}$}\\
f\!\left(\tfrac{1}{3}\right) &= \frac{1}{3}-\frac{2}{3}+\frac{3}{3} &&\justif{$3\times\frac{1}{9}=\frac{3}{9}=\frac{1}{3}$ ; $1=\frac{3}{3}$}\\
f\!\left(\tfrac{1}{3}\right) &= \res{\frac{2}{3}} &&
\end{calculs}
\cas{2)}
\begin{calculs}
g(2) &= -2^{2}+6\times 2 &&\justif{$-2^{2}=-4$ : le carré porte sur $2$ seulement}\\
g(2) &= -4+12=\res{8} &&\\[3pt]
g(-1) &= -(-1)^{2}+6\times(-1) &&\justif{$-(-1)^{2}=-1$}\\
g(-1) &= -1-6=\res{-7} &&\\[3pt]
g(7) &= -7^{2}+6\times 7 &&\\
g(7) &= -49+42=\res{-7} &&\\[3pt]
g\!\left(\tfrac{1}{2}\right) &= -\left(\frac{1}{2}\right)^{2}+6\times\frac{1}{2} &&\\
g\!\left(\tfrac{1}{2}\right) &= -\frac{1}{4}+3 &&\\
g\!\left(\tfrac{1}{2}\right) &= -\frac{1}{4}+\frac{12}{4}=\res{\frac{11}{4}} &&\justif{$3=\frac{12}{4}$}
\end{calculs}
\cas{3)}
\explique{$g(-1)=g(7)=-7$ : deux nombres différents ont \res{la même image}.}{on
l'expliquera avec la courbe (partie VII) : $-1$ et $7$ sont à égale distance de $3$}
\rem{Erreur fréquente : $(-2)^{2}=-4$. Le carré d'un nombre négatif est positif :
$(-2)^{2}=(-2)\times(-2)=4$. En revanche, $-x^{2}$ pour $x=-1$ donne $-(-1)^{2}=-1$.}
\end{exo}

\partie{Forme factorisée et racines}

\begin{exo}[De la forme factorisée à la forme développée]
\explique{On développe d'abord les deux parenthèses, puis on multiplie par le nombre
devant.}{méthode 2 du cours}
\cas{$f$}
\begin{calculs}
f(x) &= 3(x-2)(x+1) &&\justif{on recopie}\\
f(x) &= 3\left(x^{2}+x-2x-2\right) &&\justif{les deux parenthèses d'abord}\\
f(x) &= 3\left(x^{2}-x-2\right) &&\\
f(x) &= \res{3x^{2}-3x-6} &&\justif{puis on multiplie par $3$}
\end{calculs}
\cas{$g$}
\begin{calculs}
g(x) &= -(x+3)(x-4) &&\justif{on recopie}\\
g(x) &= -\left(x^{2}-4x+3x-12\right) &&\justif{les deux parenthèses d'abord}\\
g(x) &= -\left(x^{2}-x-12\right) &&\\
g(x) &= \res{-x^{2}+x+12} &&\justif{le $-$ change tous les signes}
\end{calculs}
\cas{$h$}
\begin{calculs}
h(x) &= 2(x+4)(x+1) &&\justif{on recopie}\\
h(x) &= 2\left(x^{2}+x+4x+4\right) &&\\
h(x) &= 2\left(x^{2}+5x+4\right) &&\\
h(x) &= \res{2x^{2}+10x+8} &&\justif{puis on multiplie par $2$}
\end{calculs}
\cas{$k$}
\begin{calculs}
k(x) &= -3(x-1)(x-2) &&\justif{on recopie}\\
k(x) &= -3\left(x^{2}-2x-x+2\right) &&\justif{$(-1)\times(-2)=+2$}\\
k(x) &= -3\left(x^{2}-3x+2\right) &&\\
k(x) &= \res{-3x^{2}+9x-6} &&\justif{puis on multiplie par $-3$}
\end{calculs}
\rem{Erreur fréquente : multiplier le nombre devant par les \emph{deux} parenthèses,
$3(x-2)(x+1)=(3x-6)(3x+3)$. Le $3$ ne multiplie qu'une fois.}
\end{exo}

\begin{exo}[Équations produit]
\explique{Un produit est nul si l'un de ses facteurs est nul ; un nombre comme $-3$ ou
$-4$ ne s'annule jamais.}{méthode 2 du cours}
\cas{a)}
\begin{calculs}
(x-5)(x+2) &= 0 &&\justif{on recopie}\\
x-5 &= 0 \quad\text{ou}\quad x+2=0 &&\justif{produit nul}\\
x &= 5 \quad\text{ou}\quad x=-2 &&
\end{calculs}
\explique{$\res{S=\ens{-2\,;5}}$.}{on range les solutions dans l'ordre croissant}
\cas{b)}
\begin{calculs}
-3(x+4)(2x-3) &= 0 &&\justif{on recopie}\\
x+4 &= 0 \quad\text{ou}\quad 2x-3=0 &&\justif{le facteur $-3$ ne s'annule pas}\\
x &= -4 \quad\text{ou}\quad 2x=3 &&\\
x &= -4 \quad\text{ou}\quad x=\frac{3}{2} &&\justif{on divise par $2$}
\end{calculs}
\explique{$\res{S=\ens{-4\,;\tfrac{3}{2}}}$.}{}
\cas{c)}
\begin{calculs}
-4(x+3)(x+7) &= 0 &&\justif{on recopie}\\
x+3 &= 0 \quad\text{ou}\quad x+7=0 &&\justif{le facteur $-4$ ne s'annule pas}\\
x &= -3 \quad\text{ou}\quad x=-7 &&
\end{calculs}
\explique{$\res{S=\ens{-7\,;-3}}$.}{}
\cas{d)}
\begin{calculs}
x(2x-8) &= 0 &&\justif{on recopie}\\
x &= 0 \quad\text{ou}\quad 2x-8=0 &&\justif{le premier facteur est $x$}\\
x &= 0 \quad\text{ou}\quad 2x=8 &&\\
x &= 0 \quad\text{ou}\quad x=4 &&
\end{calculs}
\explique{$\res{S=\ens{0\,;4}}$.}{}
\rem{En d), on n'oublie pas la solution $0$ : le facteur $x$ s'annule en $0$.}
\end{exo}

\begin{exo}[Lire les racines]
\explique{Dans $a(x-x_{1})(x-x_{2})$, les racines se lisent en changeant le signe :
dans $(x-3)$, la racine est $3$ ; dans $(x+4)$, elle est $-4$.}{propriété 1 du cours}
\cas{$f(x)=(x-3)(x+4)$}
\explique{Racines \res{$3$ et $-4$}.}{}
\cas{$g(x)=2(x+2)(x-6)$}
\explique{Racines \res{$-2$ et $6$}.}{le $2$ devant ne s'annule pas}
\cas{$h(x)=-3(x-5)(x+1)$}
\explique{Racines \res{$5$ et $-1$}.}{}
\cas{$k(x)=(3x-1)(x+2)$}
\explique{$3x-1=0$ donne $3x=1$, soit $x=\dfrac{1}{3}$. Racines
\res{$\dfrac{1}{3}$ et $-2$}.}{devant $x$, il y a $3$ : on ne lit pas directement}
\cas{$l(x)=4(x+3)(x+5)$}
\explique{Racines \res{$-3$ et $-5$}.}{}
\cas{$m(x)=-(x+5)(x-5)$}
\explique{Racines \res{$-5$ et $5$}.}{}
\cas{$n(x)=x(x-7)$}
\explique{Racines \res{$0$ et $7$}.}{$x=x-0$ : racine $0$}
\cas{$p(x)=2(x+1)(x+6)$}
\explique{Racines \res{$-1$ et $-6$}.}{}
\rem{Contrôle : on remplace $x$ par la racine trouvée, une parenthèse doit valoir $0$.
Pour $k$ : $3\times\frac{1}{3}-1=0$.}
\end{exo}

\begin{exo}[Les deux formes d'une même fonction]
\cas{1)}
\begin{calculs}
f(x) &= 3(x+1)(x-2) &&\justif{on recopie}\\
f(x) &= 3\left(x^{2}-2x+x-2\right) &&\justif{les deux parenthèses d'abord}\\
f(x) &= 3\left(x^{2}-x-2\right) &&\\
f(x) &= \res{3x^{2}-3x-6} &&\justif{puis on multiplie par $3$}
\end{calculs}
\cas{2)}
\begin{calculs}
f(0) &= 3\times(0+1)\times(0-2) &&\justif{avec la forme factorisée}\\
f(0) &= 3\times 1\times(-2)=-6 &&\\[3pt]
f(0) &= 3\times 0^{2}-3\times 0-6 &&\justif{avec la forme développée}\\
f(0) &= -6 &&
\end{calculs}
\explique{Les deux formes donnent $\res{f(0)=-6}$.}{c'est un contrôle du
développement}
\cas{3)}
\explique{Racines : \res{$-1$ et $2$}.}{lues sur la forme factorisée, en changeant le
signe}
\begin{calculs}
3(x+1)(x-2) &= 0 &&\justif{on résout avec la forme factorisée}\\
x+1 &= 0 \quad\text{ou}\quad x-2=0 &&\justif{produit nul}
\end{calculs}
\explique{$\res{S=\ens{-1\,;2}}$.}{les solutions de $f(x)=0$ sont les racines}
\cas{4)}
\explique{\res{Oui} : sa forme développée est $3x^{2}-3x-6$, avec $a=\res{3}$.}{le nombre
devant les parenthèses est $a$}
\rem{Chaque forme a son usage : la forme factorisée pour les racines, la forme développée
pour lire $a$, $b$ et $c$.}
\end{exo}

\partie{Le discriminant}

\begin{exo}[Calculer un discriminant]
\explique{On lit $a$, $b$ et $c$ avec leur signe, puis on écrit $\Delta=b^{2}-4ac$ en
mettant les nombres négatifs entre parenthèses.}{méthode 3 du cours}
\begin{calculs}
\Delta_{f} &= (-2)^{2}-4\times 1\times(-15) &&\justif{$a=1$, $b=-2$, $c=-15$}\\
\Delta_{f} &= 4+60=\res{64} &&\justif{$\Delta_{f}>0$}\\[3pt]
\Delta_{g} &= 2^{2}-4\times 1\times(-15) &&\justif{$a=1$, $b=2$, $c=-15$}\\
\Delta_{g} &= 4+60=\res{64} &&\justif{$\Delta_{g}>0$}\\[3pt]
\Delta_{h} &= 2^{2}-4\times 1\times 15 &&\justif{$a=1$, $b=2$, $c=15$}\\
\Delta_{h} &= 4-60=\res{-56} &&\justif{$\Delta_{h}<0$}\\[3pt]
\Delta_{k} &= 2^{2}-4\times(-1)\times(-15) &&\justif{$a=-1$, $b=2$, $c=-15$}\\
\Delta_{k} &= 4-60=\res{-56} &&\justif{$-4\times(-1)\times(-15)=-60$ : $\Delta_{k}<0$}\\[3pt]
\Delta_{l} &= (-12)^{2}-4\times 4\times 9 &&\justif{$a=4$, $b=-12$, $c=9$}\\
\Delta_{l} &= 144-144=\res{0} &&\\[3pt]
\Delta_{m} &= (-10)^{2}-4\times 25\times 1 &&\justif{$a=25$, $b=-10$, $c=1$}\\
\Delta_{m} &= 100-100=\res{0} &&\\[3pt]
\Delta_{n} &= 5^{2}-4\times(-2)\times(-3) &&\justif{$a=-2$, $b=5$, $c=-3$}\\
\Delta_{n} &= 25-24=\res{1} &&\justif{$\Delta_{n}>0$}\\[3pt]
\Delta_{p} &= 4^{2}-4\times(-3)\times(-2) &&\justif{$a=-3$, $b=4$, $c=-2$}\\
\Delta_{p} &= 16-24=\res{-8} &&\justif{$\Delta_{p}<0$}\\[3pt]
\Delta_{q} &= 7^{2}-4\times(-2)\times(-3) &&\justif{$a=-2$, $b=7$, $c=-3$}\\
\Delta_{q} &= 49-24=\res{25} &&\justif{$\Delta_{q}>0$}
\end{calculs}
\rem{$b^{2}$ est toujours positif : $(-2)^{2}=4$. Le signe de $-4ac$ dépend des signes de
$a$ et de $c$ : s'ils sont de même signe, $-4ac$ est négatif ($h$, $k$, $l$, $m$, $n$, $p$, $q$).}
\rem{Erreur fréquente : écrire $-2^{2}$ au lieu de $(-2)^{2}$ ; la calculatrice donne
alors $-4$.}
\end{exo}

\begin{exo}[Combien de racines ?]
\explique{Le signe de $\Delta$ donne le nombre de racines : deux si $\Delta>0$, une seule
si $\Delta=0$, aucune si $\Delta<0$.}{propriété 2 du cours}
\begin{calculs}
\Delta_{f} &= 5^{2}-4\times 1\times 4 &&\justif{$a=1$, $b=5$, $c=4$}\\
\Delta_{f} &= 25-16=9 &&\justif{$\Delta_{f}>0$ : \res{deux racines}}\\[3pt]
\Delta_{g} &= 4^{2}-4\times 1\times 5 &&\justif{$a=1$, $b=4$, $c=5$}\\
\Delta_{g} &= 16-20=-4 &&\justif{$\Delta_{g}<0$ : \res{aucune racine}}\\[3pt]
\Delta_{h} &= (-8)^{2}-4\times 1\times 7 &&\justif{$a=1$, $b=-8$, $c=7$}\\
\Delta_{h} &= 64-28=36 &&\justif{$\Delta_{h}>0$ : \res{deux racines}}\\[3pt]
\Delta_{k} &= 6^{2}-4\times 9\times 1 &&\justif{$a=9$, $b=6$, $c=1$}\\
\Delta_{k} &= 36-36=0 &&\justif{$\Delta_{k}=0$ : \res{une seule racine}}\\[3pt]
\Delta_{l} &= 5^{2}-4\times 2\times 3 &&\justif{$a=2$, $b=5$, $c=3$}\\
\Delta_{l} &= 25-24=1 &&\justif{$\Delta_{l}>0$ : \res{deux racines}}\\[3pt]
\Delta_{m} &= (-1)^{2}-4\times(-3)\times(-2) &&\justif{$a=-3$, $b=-1$, $c=-2$}\\
\Delta_{m} &= 1-24=-23 &&\justif{$\Delta_{m}<0$ : \res{aucune racine}}
\end{calculs}
\rem{Pour $l$, $\Delta_{l}=1$ est petit mais positif : il y a bien deux racines. Seul le
\emph{signe} de $\Delta$ compte.}
\end{exo}

\begin{exo}[Quand $b$ ou $c$ est nul]
\explique{Un terme absent a un coefficient nul : si $b=0$, $\Delta=-4ac$ ; si $c=0$,
$\Delta=b^{2}$.}{méthode 3 du cours}
\begin{calculs}
\Delta_{f} &= 0^{2}-4\times 1\times(-25) &&\justif{$a=1$, $b=0$, $c=-25$}\\
\Delta_{f} &= \res{100} &&\justif{$\Delta_{f}>0$ : \res{deux racines}}\\[3pt]
\Delta_{g} &= 0^{2}-4\times(-3)\times 12 &&\justif{$a=-3$, $b=0$, $c=12$}\\
\Delta_{g} &= \res{144} &&\justif{$\Delta_{g}>0$ : \res{deux racines}}\\[3pt]
\Delta_{h} &= (-7)^{2}-4\times 2\times 0 &&\justif{$a=2$, $b=-7$, $c=0$}\\
\Delta_{h} &= \res{49} &&\justif{$\Delta_{h}>0$ : \res{deux racines}}\\[3pt]
\Delta_{k} &= 6^{2}-4\times(-1)\times 0 &&\justif{$a=-1$, $b=6$, $c=0$}\\
\Delta_{k} &= \res{36} &&\justif{$\Delta_{k}>0$ : \res{deux racines}}\\[3pt]
\Delta_{l} &= 0^{2}-4\times(-2)\times(-5) &&\justif{$a=-2$, $b=0$, $c=-5$}\\
\Delta_{l} &= \res{-40} &&\justif{$\Delta_{l}<0$ : \res{aucune racine}}\\[3pt]
\Delta_{m} &= 0^{2}-4\times 4\times 1 &&\justif{$a=4$, $b=0$, $c=1$}\\
\Delta_{m} &= \res{-16} &&\justif{$\Delta_{m}<0$ : \res{aucune racine}}
\end{calculs}
\rem{Quand $c=0$ ($h$ et $k$), $\Delta=b^{2}$ est toujours positif (ou nul) : il y a
toujours la racine $0$, puisqu'on peut factoriser par $x$.}
\rem{Erreur fréquente pour $g$ : prendre $b=12$. Le nombre seul est $c$ ; il n'y a pas
de terme en $x$, donc $b=0$.}
\end{exo}

\partie{Racines et équations du second degré}

\begin{exo}[Deux racines]
\explique{On calcule $\Delta$ ; ici $\Delta>0$, donc deux solutions :
$x_{1}=\dfrac{-b-\sqrt{\Delta}}{2a}$ et $x_{2}=\dfrac{-b+\sqrt{\Delta}}{2a}$.}{méthode 4
du cours}
\cas{a)}
\begin{calculs}
\Delta &= (-2)^{2}-4\times 1\times(-15) &&\justif{$a=1$, $b=-2$, $c=-15$}\\
\Delta &= 4+60=64 &&\justif{$\Delta>0$ : deux solutions ; $\sqrt{64}=8$}\\
x_{1} &= \frac{2-8}{2}=-3 &&\justif{$-b=2$ et $2a=2$}\\
x_{2} &= \frac{2+8}{2}=5 &&
\end{calculs}
\explique{$\res{S=\ens{-3\,;5}}$.}{contrôle : $(-3)^{2}-2\times(-3)-15=0$}
\cas{b)}
\begin{calculs}
\Delta &= 3^{2}-4\times 2\times(-2) &&\justif{$a=2$, $b=3$, $c=-2$}\\
\Delta &= 9+16=25 &&\justif{$\sqrt{25}=5$}\\
x_{1} &= \frac{-3-5}{4}=-2 &&\justif{$2a=4$}\\
x_{2} &= \frac{-3+5}{4}=\frac{1}{2} &&\justif{$\frac{2}{4}=\frac{1}{2}$}
\end{calculs}
\explique{$\res{S=\ens{-2\,;\tfrac{1}{2}}}$.}{}
\cas{c)}
\begin{calculs}
\Delta &= (-8)^{2}-4\times 3\times 5 &&\justif{$a=3$, $b=-8$, $c=5$}\\
\Delta &= 64-60=4 &&\justif{$\sqrt{4}=2$}\\
x_{1} &= \frac{8-2}{6}=1 &&\justif{$-b=8$ et $2a=6$}\\
x_{2} &= \frac{8+2}{6}=\frac{5}{3} &&\justif{$\frac{10}{6}=\frac{5}{3}$}
\end{calculs}
\explique{$\res{S=\ens{1\,;\tfrac{5}{3}}}$.}{}
\cas{d)}
\begin{calculs}
\Delta &= 7^{2}-4\times(-1)\times(-12) &&\justif{$a=-1$, $b=7$, $c=-12$}\\
\Delta &= 49-48=1 &&\justif{$\sqrt{1}=1$}\\
x_{1} &= \frac{-7-1}{-2}=4 &&\justif{$2a=-2$ : on divise par un négatif}\\
x_{2} &= \frac{-7+1}{-2}=3 &&
\end{calculs}
\explique{$\res{S=\ens{3\,;4}}$.}{on range les solutions dans l'ordre croissant}
\rem{Erreur fréquente : oublier le signe de $-b$. Si $b=-2$, alors $-b=2$.}
\rem{Avec $a<0$ (en d), le dénominateur $2a$ est négatif : $x_{1}$ est alors la plus
grande solution.}
\end{exo}

\begin{exo}[Les trois cas]
\explique{Le signe de $\Delta$ dit le cas : deux solutions, une seule
$x_{0}=\dfrac{-b}{2a}$, ou aucune.}{théorème 1 du cours}
\cas{a)}
\begin{calculs}
\Delta &= (-20)^{2}-4\times 4\times 25 &&\justif{$a=4$, $b=-20$, $c=25$}\\
\Delta &= 400-400=0 &&\justif{$\Delta=0$ : une seule solution}\\
x_{0} &= \frac{20}{8}=\frac{5}{2} &&\justif{$x_{0}=\frac{-b}{2a}$, avec $2a=8$}
\end{calculs}
\explique{$\res{S=\ens{\tfrac{5}{2}}}$.}{}
\cas{b)}
\begin{calculs}
\Delta &= 2^{2}-4\times 1\times 4 &&\justif{$a=1$, $b=2$, $c=4$}\\
\Delta &= 4-16=-12 &&\justif{$\Delta<0$}
\end{calculs}
\explique{$\res{S=\varnothing}$.}{aucune solution}
\cas{c)}
\begin{calculs}
\Delta &= 2^{2}-4\times(-3)\times(-5) &&\justif{$a=-3$, $b=2$, $c=-5$}\\
\Delta &= 4-60=-56 &&\justif{$\Delta<0$}
\end{calculs}
\explique{$\res{S=\varnothing}$.}{aucune solution}
\cas{d)}
\begin{calculs}
\Delta &= 8^{2}-4\times 1\times 16 &&\justif{$a=1$, $b=8$, $c=16$}\\
\Delta &= 64-64=0 &&\justif{$\Delta=0$ : une seule solution}\\
x_{0} &= \frac{-8}{2}=-4 &&\justif{$x_{0}=\frac{-b}{2a}$}
\end{calculs}
\explique{$\res{S=\ens{-4}}$.}{}
\cas{e)}
\begin{calculs}
\Delta &= 1^{2}-4\times 6\times(-2) &&\justif{$a=6$, $b=1$, $c=-2$}\\
\Delta &= 1+48=49 &&\justif{$\Delta>0$ : deux solutions ; $\sqrt{49}=7$}\\
x_{1} &= \frac{-1-7}{12}=-\frac{2}{3} &&\justif{$\frac{-8}{12}=-\frac{2}{3}$}\\
x_{2} &= \frac{-1+7}{12}=\frac{1}{2} &&\justif{$\frac{6}{12}=\frac{1}{2}$}
\end{calculs}
\explique{$\res{S=\ens{-\tfrac{2}{3}\,;\tfrac{1}{2}}}$.}{}
\cas{f)}
\begin{calculs}
\Delta &= 5^{2}-4\times(-2)\times(-2) &&\justif{$a=-2$, $b=5$, $c=-2$}\\
\Delta &= 25-16=9 &&\justif{$\Delta>0$ : deux solutions ; $\sqrt{9}=3$}\\
x_{1} &= \frac{-5-3}{-4}=2 &&\justif{$2a=-4$}\\
x_{2} &= \frac{-5+3}{-4}=\frac{1}{2} &&
\end{calculs}
\explique{$\res{S=\ens{\tfrac{1}{2}\,;2}}$.}{}
\rem{En a) et d), $\Delta=0$ annonce un carré : $4x^{2}-20x+25=(2x-5)^{2}$ et
$x^{2}+8x+16=(x+4)^{2}$.}
\end{exo}

\begin{exo}[Sans discriminant]
\explique{Pas de terme en $x$ : on isole $x^{2}$, ou on reconnaît une différence de deux
carrés. Pas de nombre seul : on factorise par $x$.}{méthode 4 du cours, cas e)}
\cas{a)}
\begin{calculs}
x^{2}-7 &= 0 &&\justif{on recopie}\\
x^{2} &= 7 &&\justif{on ajoute $7$}\\
x &= -\sqrt{7} \quad\text{ou}\quad x=\sqrt{7} &&\justif{$7>0$ : deux solutions}
\end{calculs}
\explique{$\res{S=\ens{-\sqrt{7}\,;\sqrt{7}}}$.}{}
\cas{b)}
\begin{calculs}
2x^{2}-9x &= 0 &&\justif{on recopie}\\
x(2x-9) &= 0 &&\justif{facteur commun $x$}\\
x &= 0 \quad\text{ou}\quad 2x-9=0 &&\justif{produit nul}\\
x &= 0 \quad\text{ou}\quad x=\frac{9}{2} &&
\end{calculs}
\explique{$\res{S=\ens{0\,;\tfrac{9}{2}}}$.}{}
\cas{c)}
\begin{calculs}
x^{2}-64 &= 0 &&\justif{on recopie}\\
(x-8)(x+8) &= 0 &&\justif{$a^{2}-b^{2}=(a-b)(a+b)$, avec $64=8^{2}$}\\
x-8 &= 0 \quad\text{ou}\quad x+8=0 &&\justif{produit nul}
\end{calculs}
\explique{$\res{S=\ens{-8\,;8}}$.}{}
\cas{d)}
\begin{calculs}
9x^{2}-4 &= 0 &&\justif{on recopie}\\
(3x-2)(3x+2) &= 0 &&\justif{$9x^{2}=(3x)^{2}$ et $4=2^{2}$}\\
3x-2 &= 0 \quad\text{ou}\quad 3x+2=0 &&\justif{produit nul}\\
x &= \frac{2}{3} \quad\text{ou}\quad x=-\frac{2}{3} &&
\end{calculs}
\explique{$\res{S=\ens{-\tfrac{2}{3}\,;\tfrac{2}{3}}}$.}{}
\cas{e)}
\begin{calculs}
3x^{2}+12x &= 0 &&\justif{on recopie}\\
3x(x+4) &= 0 &&\justif{facteur commun $3x$}\\
3x &= 0 \quad\text{ou}\quad x+4=0 &&\justif{produit nul}
\end{calculs}
\explique{$\res{S=\ens{-4\,;0}}$.}{}
\cas{f)}
\begin{calculs}
x^{2}+9 &= 0 &&\justif{on recopie}\\
x^{2} &= -9 &&\justif{on retranche $9$}
\end{calculs}
\explique{$\res{S=\varnothing}$.}{un carré n'est jamais négatif}
\rem{Erreur fréquente en b) : diviser par $x$ et « perdre » la solution $0$. On
factorise, on ne divise pas par $x$.}
\rem{En a), $x^{2}=7$ a deux solutions : on n'oublie pas $-\sqrt{7}$.}
\end{exo}

\begin{exo}[Tout ramener d'un côté]
\explique{On développe s'il le faut, puis on met tout dans le membre de gauche pour se
ramener à $ax^{2}+bx+c=0$.}{méthode 4 du cours, cas d)}
\cas{a)}
\begin{calculs}
3x^{2}+2x-1 &= 2x^{2}+7 &&\justif{on recopie}\\
x^{2}+2x-8 &= 0 &&\justif{on retranche $2x^{2}$ et $7$}
\end{calculs}
\begin{calculs}
\Delta &= 2^{2}-4\times 1\times(-8) &&\\
\Delta &= 4+32=36 &&\justif{$\sqrt{36}=6$}\\
x_{1} &= \frac{-2-6}{2}=-4 &&\\
x_{2} &= \frac{-2+6}{2}=2 &&
\end{calculs}
\explique{$\res{S=\ens{-4\,;2}}$.}{}
\cas{b)}
\begin{calculs}
x(x+3) &= -(2x+6) &&\justif{on recopie}\\
x^{2}+3x &= -2x-6 &&\justif{on développe les deux membres}\\
x^{2}+5x+6 &= 0 &&\justif{on ajoute $2x$ et $6$}
\end{calculs}
\begin{calculs}
\Delta &= 5^{2}-4\times 1\times 6 &&\\
\Delta &= 25-24=1 &&\\
x_{1} &= \frac{-5-1}{2}=-3 &&\\
x_{2} &= \frac{-5+1}{2}=-2 &&
\end{calculs}
\explique{$\res{S=\ens{-3\,;-2}}$.}{}
\cas{c)}
\begin{calculs}
x^{2}-3x+7 &= 5 &&\justif{on recopie}\\
x^{2}-3x+2 &= 0 &&\justif{on retranche $5$}
\end{calculs}
\begin{calculs}
\Delta &= (-3)^{2}-4\times 1\times 2 &&\\
\Delta &= 9-8=1 &&\\
x_{1} &= \frac{3-1}{2}=1 &&\\
x_{2} &= \frac{3+1}{2}=2 &&
\end{calculs}
\explique{$\res{S=\ens{1\,;2}}$.}{}
\cas{d)}
\begin{calculs}
(x+2)(x-3) &= 6 &&\justif{on recopie}\\
x^{2}-3x+2x-6 &= 6 &&\justif{on développe}\\
x^{2}-x-12 &= 0 &&\justif{on retranche $6$}
\end{calculs}
\begin{calculs}
\Delta &= (-1)^{2}-4\times 1\times(-12) &&\\
\Delta &= 1+48=49 &&\justif{$\sqrt{49}=7$}\\
x_{1} &= \frac{1-7}{2}=-3 &&\\
x_{2} &= \frac{1+7}{2}=4 &&
\end{calculs}
\explique{$\res{S=\ens{-3\,;4}}$.}{contrôle : $(-3+2)\times(-3-3)=6$}
\rem{En d), on ne résout pas « $x+2=6$ ou $x-3=6$ » : la règle du produit nul ne marche
que si le produit vaut $0$.}
\end{exo}

\begin{exo}[Racines d'une fonction]
\explique{Les racines de $f$ sont les solutions de $f(x)=0$ : on résout avec
$\Delta$.}{méthode 4 du cours}
\cas{a)}
\begin{calculs}
\Delta &= (-9)^{2}-4\times 1\times 20 &&\justif{$a=1$, $b=-9$, $c=20$}\\
\Delta &= 81-80=1 &&\\
x_{1} &= \frac{9-1}{2}=4 &&\\
x_{2} &= \frac{9+1}{2}=5 &&
\end{calculs}
\explique{$f$ a deux racines : \res{$4$ et $5$}.}{}
\cas{b)}
\begin{calculs}
\Delta &= 3^{2}-4\times(-3)\times 18 &&\justif{$a=-3$, $b=3$, $c=18$}\\
\Delta &= 9+216=225 &&\justif{$\sqrt{225}=15$}\\
x_{1} &= \frac{-3-15}{-6}=3 &&\justif{$2a=-6$}\\
x_{2} &= \frac{-3+15}{-6}=-2 &&
\end{calculs}
\explique{$g$ a deux racines : \res{$-2$ et $3$}.}{}
\cas{c)}
\begin{calculs}
\Delta &= 4^{2}-4\times 1\times 7 &&\justif{$a=1$, $b=4$, $c=7$}\\
\Delta &= 16-28=-12 &&\justif{$\Delta<0$}
\end{calculs}
\explique{$h$ n'a \res{aucune racine}.}{}
\cas{d)}
\begin{calculs}
\Delta &= (-6)^{2}-4\times 9\times 1 &&\justif{$a=9$, $b=-6$, $c=1$}\\
\Delta &= 36-36=0 &&\justif{une seule racine}\\
x_{0} &= \frac{6}{18}=\frac{1}{3} &&\justif{$x_{0}=\frac{-b}{2a}$}
\end{calculs}
\explique{$k$ a une seule racine : \res{$\dfrac{1}{3}$}.}{}
\rem{En b), on peut d'abord diviser l'équation $-3x^{2}+3x+18=0$ par $-3$ :
$x^{2}-x-6=0$ a les mêmes solutions, avec de plus petits nombres.}
\end{exo}

\partie{Factoriser un trinôme}

\begin{exo}[Les racines sont données]
\explique{Avec deux racines, $P(x)=a(x-x_{1})(x-x_{2})$ ; avec une racine double,
$P(x)=a(x-x_{0})^{2}$. Le coefficient $a$ se lit devant $x^{2}$.}{théorème 2 du cours}
\cas{a)}
\explique{$a=1$, racines $-3$ et $2$ : $x^{2}+x-6=\res{(x+3)(x-2)}$.}{$x-(-3)=x+3$}
\cas{b)}
\explique{$a=3$, racines $-3$ et $2$ : $3x^{2}+3x-18=\res{3(x+3)(x-2)}$.}{mêmes racines
qu'en a), mais $a=3$}
\cas{c)}
\explique{$a=-1$, racines $1$ et $4$ : $-x^{2}+5x-4=\res{-(x-1)(x-4)}$.}{$-1$ devant
s'écrit « $-$ »}
\cas{d)}
\explique{$a=2$, racine double $3$ : $2x^{2}-12x+18=\res{2(x-3)^{2}}$.}{}
\explique{Contrôle en développant b) :}{}
\begin{calculs}
3(x+3)(x-2) &= 3\left(x^{2}-2x+3x-6\right) &&\\
3(x+3)(x-2) &= 3\left(x^{2}+x-6\right) &&\\
3(x+3)(x-2) &= 3x^{2}+3x-18 &&\justif{on retrouve le trinôme de départ}
\end{calculs}
\rem{Erreur fréquente : oublier $a$, et écrire $(x+3)(x-2)$ en b). Contrôle rapide sur le
nombre seul : $3\times 3\times(-2)=-18$, c'est bien $c$.}
\end{exo}

\begin{exo}[Factoriser]
\explique{On calcule $\Delta$, puis les racines, et on écrit $a(x-x_{1})(x-x_{2})$ en
n'oubliant pas $a$.}{méthode 5 du cours}
\cas{$f$}
\begin{calculs}
\Delta &= (-2)^{2}-4\times 1\times(-8) &&\justif{$a=1$, $b=-2$, $c=-8$}\\
\Delta &= 4+32=36 &&\justif{$\sqrt{36}=6$}\\
x_{1} &= \frac{2-6}{2}=-2 &&\\
x_{2} &= \frac{2+6}{2}=4 &&
\end{calculs}
\explique{$f(x)=\res{(x+2)(x-4)}$.}{contrôle : $2\times(-4)=-8$, c'est bien $c$}
\cas{$g$}
\begin{calculs}
\Delta &= (-1)^{2}-4\times 2\times(-1) &&\justif{$a=2$, $b=-1$, $c=-1$}\\
\Delta &= 1+8=9 &&\justif{$\sqrt{9}=3$}\\
x_{1} &= \frac{1-3}{4}=-\frac{1}{2} &&\\
x_{2} &= \frac{1+3}{4}=1 &&
\end{calculs}
\explique{$g(x)=\res{2\left(x+\tfrac{1}{2}\right)(x-1)}$.}{contrôle :
$2\times\frac{1}{2}\times(-1)=-1$}
\cas{$h$}
\begin{calculs}
\Delta &= (-2)^{2}-4\times(-1)\times 3 &&\justif{$a=-1$, $b=-2$, $c=3$}\\
\Delta &= 4+12=16 &&\justif{$\sqrt{16}=4$}\\
x_{1} &= \frac{2-4}{-2}=1 &&\justif{$2a=-2$}\\
x_{2} &= \frac{2+4}{-2}=-3 &&
\end{calculs}
\explique{$h(x)=\res{-(x-1)(x+3)}$.}{contrôle : $-1\times(-1)\times 3=3$}
\cas{$k$}
\begin{calculs}
\Delta &= (-7)^{2}-4\times 3\times 2 &&\justif{$a=3$, $b=-7$, $c=2$}\\
\Delta &= 49-24=25 &&\justif{$\sqrt{25}=5$}\\
x_{1} &= \frac{7-5}{6}=\frac{1}{3} &&\\
x_{2} &= \frac{7+5}{6}=2 &&
\end{calculs}
\explique{$k(x)=\res{3\left(x-\tfrac{1}{3}\right)(x-2)}$.}{contrôle :
$3\times\left(-\frac{1}{3}\right)\times(-2)=2$}
\rem{Pour $g$ et $k$, on peut faire entrer $a$ dans la première parenthèse :
$2\left(x+\frac{1}{2}\right)=2x+1$ et $3\left(x-\frac{1}{3}\right)=3x-1$, d'où
$g(x)=(2x+1)(x-1)$ et $k(x)=(3x-1)(x-2)$. Ce n'est pas exigé.}
\end{exo}

\begin{exo}[Factoriser si c'est possible]
\explique{Si $\Delta>0$ : $a(x-x_{1})(x-x_{2})$ ; si $\Delta=0$ : $a(x-x_{0})^{2}$ ; si
$\Delta<0$ : pas de factorisation.}{théorème 2 du cours}
\cas{$f$}
\begin{calculs}
\Delta &= 4^{2}-4\times 1\times 4 &&\justif{$a=1$, $b=4$, $c=4$}\\
\Delta &= 16-16=0 &&\justif{une racine double}\\
x_{0} &= \frac{-4}{2}=-2 &&\justif{$x_{0}=\frac{-b}{2a}$}
\end{calculs}
\explique{$f(x)=\res{(x+2)^{2}}$.}{$\Delta=0$ : un carré}
\cas{$g$}
\begin{calculs}
\Delta &= 0^{2}-4\times 1\times 9 &&\justif{$a=1$, $b=0$, $c=9$}\\
\Delta &= -36 &&\justif{$\Delta<0$}
\end{calculs}
\explique{$g$ \res{ne se factorise pas}.}{$x^{2}+9\ge 9$ : jamais nul}
\cas{$h$}
\begin{calculs}
\Delta &= 0^{2}-4\times 1\times(-9) &&\justif{$a=1$, $b=0$, $c=-9$}\\
\Delta &= 36 &&\justif{$\sqrt{36}=6$}\\
x_{1} &= \frac{0-6}{2}=-3 &&\\
x_{2} &= \frac{0+6}{2}=3 &&
\end{calculs}
\explique{$h(x)=\res{(x+3)(x-3)}$.}{plus vite : $x^{2}-3^{2}=(x-3)(x+3)$}
\cas{$k$}
\begin{calculs}
\Delta &= 1^{2}-4\times 6\times(-2) &&\justif{$a=6$, $b=1$, $c=-2$}\\
\Delta &= 1+48=49 &&\justif{$\sqrt{49}=7$}\\
x_{1} &= \frac{-1-7}{12}=-\frac{2}{3} &&\\
x_{2} &= \frac{-1+7}{12}=\frac{1}{2} &&
\end{calculs}
\explique{$k(x)=\res{6\left(x+\tfrac{2}{3}\right)\left(x-\tfrac{1}{2}\right)}$.}{contrôle :
$6\times\frac{2}{3}\times\left(-\frac{1}{2}\right)=-2$}
\cas{$l$}
\begin{calculs}
\Delta &= 1^{2}-4\times 3\times 2 &&\justif{$a=3$, $b=1$, $c=2$}\\
\Delta &= 1-24=-23 &&\justif{$\Delta<0$}
\end{calculs}
\explique{$l$ \res{ne se factorise pas}.}{}
\cas{$m$}
\begin{calculs}
\Delta &= 6^{2}-4\times(-3)\times(-3) &&\justif{$a=-3$, $b=6$, $c=-3$}\\
\Delta &= 36-36=0 &&\justif{une racine double}\\
x_{0} &= \frac{-6}{-6}=1 &&\justif{$x_{0}=\frac{-b}{2a}$}
\end{calculs}
\explique{$m(x)=\res{-3(x-1)^{2}}$.}{on n'oublie pas $a=-3$ devant le carré}
\rem{« Ne se factorise pas » se justifie toujours par le calcul : $\Delta<0$.}
\end{exo}

\begin{exo}[Racine évidente]
\explique{On vérifie que l'image du nombre vaut $0$. On écrit alors
$a(x-x_{1})(x-x_{2})$ avec $x_{2}$ inconnu, et on trouve $x_{2}$ en remplaçant $x$ par
$0$ dans les deux formes.}{remarque « racine évidente » de la partie V}
\cas{a)}
\begin{calculs}
f(1) &= 1^{2}-6\times 1+5 &&\\
f(1) &= 1-6+5=0 &&\justif{$1$ est bien une racine}
\end{calculs}
\explique{$a=1$, donc $f(x)=(x-1)(x-x_{2})$. En $x=0$, les deux formes donnent
$f(0)$ :}{$f(0)=5$, le nombre seul}
\begin{calculs}
(0-1)\times(0-x_{2}) &= 5 &&\justif{forme factorisée en $0$ ; forme développée : $5$}\\
(-1)\times(-x_{2}) &= 5 &&\\
x_{2} &= 5 &&
\end{calculs}
\explique{$f(x)=\res{(x-1)(x-5)}$.}{contrôle : $f(5)=25-30+5=0$}
\cas{b)}
\begin{calculs}
g(1) &= 2\times 1^{2}+5\times 1-7 &&\\
g(1) &= 2+5-7=0 &&\justif{$1$ est bien une racine}
\end{calculs}
\explique{$a=2$, donc $g(x)=2(x-1)(x-x_{2})$. En $x=0$ :}{$g(0)=-7$}
\begin{calculs}
2\times(0-1)\times(0-x_{2}) &= -7 &&\\
2x_{2} &= -7 &&\justif{$2\times(-1)\times(-x_{2})=2x_{2}$}\\
x_{2} &= -\frac{7}{2} &&
\end{calculs}
\explique{$g(x)=\res{2(x-1)\left(x+\tfrac{7}{2}\right)}$.}{ou $(x-1)(2x+7)$}
\cas{c)}
\begin{calculs}
h(-1) &= -2\times(-1)^{2}-5\times(-1)-3 &&\\
h(-1) &= -2+5-3=0 &&\justif{$-1$ est bien une racine}
\end{calculs}
\explique{$a=-2$, donc $h(x)=-2(x+1)(x-x_{2})$. En $x=0$ :}{$h(0)=-3$}
\begin{calculs}
-2\times(0+1)\times(0-x_{2}) &= -3 &&\\
2x_{2} &= -3 &&\justif{$-2\times 1\times(-x_{2})=2x_{2}$}\\
x_{2} &= -\frac{3}{2} &&
\end{calculs}
\explique{$h(x)=\res{-2(x+1)\left(x+\tfrac{3}{2}\right)}$.}{ou $-(x+1)(2x+3)$}
\rem{Le calcul de $\Delta$ donne le même résultat, en plus long. Pour $g$ :
$\Delta=25+56=81$, racines $\frac{-5-9}{4}=-\frac{7}{2}$ et $\frac{-5+9}{4}=1$.}
\end{exo}

\partie{Représentation graphique}

\begin{exo}[Tableau de valeurs et tracé]
\cas{1)}
\explique{$a=\res{1>0}$ : la parabole est \res{un creux}.}{propriété 3 du cours :
tournée vers le haut}
\cas{2)}
\begin{calculs}
\Delta &= 1^{2}-4\times 1\times(-2) &&\justif{$a=1$, $b=1$, $c=-2$}\\
\Delta &= 1+8=9 &&\justif{$\Delta>0$ : deux racines ; $\sqrt{9}=3$}\\
x_{1} &= \frac{-1-3}{2}=-2 &&\\
x_{2} &= \frac{-1+3}{2}=1 &&
\end{calculs}
\explique{Les racines de $f$ sont \res{$-2$ et $1$}.}{la courbe coupera l'axe des
abscisses en ces deux points}
\cas{3)}
\begin{calculs}
f(-3) &= 9-3-2=4 &&\justif{$(-3)^{2}=9$}\\
f(-1) &= 1-1-2=-2 &&\justif{$(-1)^{2}=1$}\\
f(0) &= -2 &&\justif{le nombre seul}\\
f(2) &= 4+2-2=4 &&
\end{calculs}
\explique{$f(-2)=0$ et $f(1)=0$ : ce sont les racines.}{pas besoin de les recalculer}
\begin{center}
\renewcommand{\arraystretch}{1.25}
\begin{tabular}{|c|c|c|c|c|c|c|}
\hline
$x$ & $-3$ & $-2$ & $-1$ & $0$ & $1$ & $2$\\
\hline
$f(x)$ & \res{$4$} & \res{$0$} & \res{$-2$} & \res{$-2$} & \res{$0$} & \res{$4$}\\
\hline
\end{tabular}
\end{center}
\cas{4)}
\begin{minipage}[c]{0.58\linewidth}\raggedright
\explique{On place les six points du tableau, puis on les relie à main levée par une
courbe arrondie, sans segments.}{méthode 6 du cours}
\explique{Le bas de la courbe est entre $x=-1$ et $x=0$ : $f(-1)=f(0)$.}{les valeurs se
répondent : $f(-3)=f(2)$, $f(-2)=f(1)$}
\end{minipage}\hfill
\begin{minipage}[c]{0.38\linewidth}\centering
\begin{tikzpicture}[x=0.5cm,y=0.5cm]
  \repere{-4}{-3}{3}{5}
  \gradx{-3,-2,-1,1,2} \grady{-2,-1,1,2,3,4}
  \draw[coulCourbe,line width=1.1pt,domain=-3:2,samples=60] plot (\x,{\x*\x+\x-2});
  \foreach \p in {(-3,4),(-2,0),(-1,-2),(0,-2),(1,0),(2,4)} \path[coulCourbe] \p node[croix]{};
  \node[coulCourbe,font=\footnotesize,right] at (2,4) {$\Cf$};
\end{tikzpicture}
\end{minipage}
\rem{Erreur fréquente : relier les points par des segments, ou mettre la pointe du creux
sur un point du tableau. Le bas de la parabole est à $x=-\frac{1}{2}$, entre deux points.}
\end{exo}

\begin{exo}[Un deuxième tracé]
\explique{$a=\res{-1<0}$ : la parabole est \res{une bosse}.}{propriété 3 du cours}
\begin{calculs}
\Delta &= 4^{2}-4\times(-1)\times(-3) &&\justif{$a=-1$, $b=4$, $c=-3$}\\
\Delta &= 16-12=4 &&\justif{$\sqrt{4}=2$}\\
x_{1} &= \frac{-4-2}{-2}=3 &&\justif{$2a=-2$}\\
x_{2} &= \frac{-4+2}{-2}=1 &&
\end{calculs}
\explique{Racines : \res{$1$ et $3$}.}{}
\begin{calculs}
g(0) &= -3 &&\justif{le nombre seul}\\
g(2) &= -4+8-3=1 &&\justif{$-2^{2}=-4$}\\
g(4) &= -16+16-3=-3 &&\justif{$-4^{2}=-16$}
\end{calculs}
\begin{minipage}[c]{0.58\linewidth}\raggedright
\begin{center}
\renewcommand{\arraystretch}{1.25}
\begin{tabular}{|c|c|c|c|c|c|}
\hline
$x$ & $0$ & $1$ & $2$ & $3$ & $4$\\
\hline
$g(x)$ & \res{$-3$} & \res{$0$} & \res{$1$} & \res{$0$} & \res{$-3$}\\
\hline
\end{tabular}
\end{center}
\explique{On place les points et on relie par une courbe arrondie ; le haut de la bosse
est en $(2\,;1)$.}{symétrie : $g(0)=g(4)$ et $g(1)=g(3)$}
\end{minipage}\hfill
\begin{minipage}[c]{0.38\linewidth}\centering
\begin{tikzpicture}[x=0.5cm,y=0.5cm]
  \repere{-1}{-4}{5}{2}
  \gradx{1,2,3,4} \grady{-3,-2,-1,1}
  \draw[coulCourbe,line width=1.1pt,domain=0:4,samples=60] plot (\x,{-\x*\x+4*\x-3});
  \foreach \p in {(0,-3),(1,0),(2,1),(3,0),(4,-3)} \path[coulCourbe] \p node[croix]{};
  \node[coulCourbe,font=\footnotesize,right] at (4,-3) {$\Cg$};
\end{tikzpicture}
\end{minipage}
\rem{Erreur fréquente : $g(2)=-(-4)+8-3$. Le carré porte sur $2$ seulement :
$-2^{2}=-4$.}
\end{exo}

\begin{exo}[Associer une courbe à une fonction]
\cas{1)}
\explique{$f(x)=x^{2}-1$ : $a=1>0$ ; $\Delta=0^{2}-4\times 1\times(-1)=4>0$ : \res{deux
racines}.}{$x^{2}-1=(x-1)(x+1)$ : racines $-1$ et $1$}
\explique{$g(x)=-x^{2}+3$ : $a=-1<0$ ; $\Delta=0^{2}-4\times(-1)\times 3=12>0$ :
\res{deux racines}.}{}
\explique{$h(x)=x^{2}+1$ : $a=1>0$ ; $\Delta=0^{2}-4\times 1\times 1=-4<0$ : \res{aucune
racine}.}{}
\explique{$k(x)=(x+1)^{2}=x^{2}+2x+1$ : $a=1>0$ ; $\Delta=2^{2}-4\times 1\times 1=0$ :
\res{une seule racine}, $-1$.}{on développe pour lire $a$, $b$ et $c$}
\cas{2)}
\explique{Seule $g$ a $a<0$ : c'est la seule bosse, \res{courbe 1}.}{signe de $a$}
\explique{$k$ a une seule racine : sa courbe touche l'axe des abscisses en un point,
\res{courbe 2}.}{en $x=-1$}
\explique{$f$ a deux racines : sa courbe coupe l'axe en deux points, \res{courbe 3}.}{en
$-1$ et $1$}
\explique{$h$ n'a pas de racine : sa courbe ne coupe pas l'axe, \res{courbe 4}.}{au-dessus
de l'axe}
\rem{Contrôle par une image : $f(0)=-1$, et la courbe 3 passe bien par $(0\,;-1)$ ;
$h(0)=1$, et la courbe 4 par $(0\,;1)$.}
\end{exo}

\begin{exo}[Sans tracer]
\explique{La parabole coupe l'axe des abscisses en deux points si $\Delta>0$, le touche en
un point si $\Delta=0$, ne le coupe pas si $\Delta<0$.}{propriété 4 du cours}
\begin{calculs}
\Delta_{f} &= (-8)^{2}-4\times 1\times 16=0 &&\justif{\res{un point}}\\
\Delta_{g} &= 1^{2}-4\times 3\times 4=-47 &&\justif{\res{aucun point}}\\
\Delta_{h} &= 2^{2}-4\times(-1)\times 8=36 &&\justif{\res{deux points}}\\
\Delta_{k} &= 0^{2}-4\times 2\times(-18)=144 &&\justif{\res{deux points}}\\
\Delta_{l} &= 6^{2}-4\times(-9)\times(-1)=0 &&\justif{\res{un point}}\\
\Delta_{m} &= 4^{2}-4\times 1\times 6=-8 &&\justif{\res{aucun point}}
\end{calculs}
\rem{Pour $k$, $b=0$ : il n'y a pas de terme en $x$.}
\end{exo}

\partie{Sommet et variations}

\begin{exo}[Coordonnées du sommet]
\explique{$\alpha=-\dfrac{b}{2a}$, puis $\beta=f(\alpha)$ : on remplace $x$ par
$\alpha$.}{théorème 3 du cours}
\cas{a)}
\begin{calculs}
\alpha &= -\frac{-6}{2\times 1} &&\justif{$a=1$ et $b=-6$}\\
\alpha &= 3 &&\\
\beta &= 3^{2}-6\times 3+2 &&\justif{$\beta=f(3)$}\\
\beta &= 9-18+2=-7 &&
\end{calculs}
\explique{$\res{S(3\,;-7)}$.}{}
\cas{b)}
\begin{calculs}
\alpha &= -\frac{4}{2\times(-1)} &&\justif{$a=-1$ et $b=4$}\\
\alpha &= 2 &&\justif{$-\frac{4}{-2}=2$}\\
\beta &= -2^{2}+4\times 2+3 &&\justif{$\beta=g(2)$}\\
\beta &= -4+8+3=7 &&
\end{calculs}
\explique{$\res{S(2\,;7)}$.}{}
\cas{c)}
\begin{calculs}
\alpha &= -\frac{8}{2\times 2} &&\justif{$a=2$ et $b=8$}\\
\alpha &= -2 &&\\
\beta &= 2\times(-2)^{2}+8\times(-2)+5 &&\justif{$\beta=h(-2)$}\\
\beta &= 8-16+5=-3 &&
\end{calculs}
\explique{$\res{S(-2\,;-3)}$.}{}
\cas{d)}
\begin{calculs}
\alpha &= -\frac{12}{2\times(-2)} &&\justif{$a=-2$ et $b=12$}\\
\alpha &= 3 &&\justif{$-\frac{12}{-4}=3$}\\
\beta &= -2\times 3^{2}+12\times 3-7 &&\justif{$\beta=k(3)$}\\
\beta &= -18+36-7=11 &&
\end{calculs}
\explique{$\res{S(3\,;11)}$.}{}
\rem{Erreur fréquente : oublier le « $-$ » de $-\frac{b}{2a}$. En a), $b=-6$ donne
$\alpha=-\frac{-6}{2}=3$, et non $-3$.}
\end{exo}

\begin{exo}[Tableau de variations]
\explique{On calcule le sommet ; si $a>0$, la fonction décroît puis croît (minimum
$\beta$) ; si $a<0$, elle croît puis décroît (maximum $\beta$).}{méthode 7 du cours}
\cas{a)}
\begin{calculs}
\alpha &= -\frac{-4}{2\times 1}=2 &&\justif{$a=1$ et $b=-4$}\\
\beta &= 2^{2}-4\times 2+3=-1 &&\justif{$\beta=f(2)$}
\end{calculs}
\tabvar{f(x)}{$-\infty$,$2$,$+\infty$}{+/ ,-/$-1$,+/ }
\explique{$f$ admet un \res{minimum, $-1$}, atteint en $x=2$.}{$a=1>0$ : un creux}
\cas{b)}
\begin{calculs}
\alpha &= -\frac{6}{2\times(-3)}=1 &&\justif{$a=-3$ et $b=6$}\\
\beta &= -3\times 1^{2}+6\times 1+9=12 &&\justif{$\beta=g(1)$}
\end{calculs}
\tabvar{g(x)}{$-\infty$,$1$,$+\infty$}{-/ ,+/$12$,-/ }
\explique{$g$ admet un \res{maximum, $12$}, atteint en $x=1$.}{$a=-3<0$ : une bosse}
\cas{c)}
\begin{calculs}
\alpha &= -\frac{8}{2\times 2}=-2 &&\justif{$a=2$, $b=8$ ($c=0$)}\\
\beta &= 2\times(-2)^{2}+8\times(-2)=-8 &&\justif{$8-16=-8$}
\end{calculs}
\tabvar{h(x)}{$-\infty$,$-2$,$+\infty$}{+/ ,-/$-8$,+/ }
\explique{$h$ admet un \res{minimum, $-8$}, atteint en $x=-2$.}{$a=2>0$}
\cas{d)}
\begin{calculs}
\alpha &= -\frac{-6}{2\times(-1)}=-3 &&\justif{$a=-1$ et $b=-6$ : $-\frac{-6}{-2}=-3$}\\
\beta &= -(-3)^{2}-6\times(-3)-5=4 &&\justif{$-9+18-5=4$}
\end{calculs}
\tabvar{k(x)}{$-\infty$,$-3$,$+\infty$}{-/ ,+/$4$,-/ }
\explique{$k$ admet un \res{maximum, $4$}, atteint en $x=-3$.}{$a=-1<0$}
\rem{En d), deux signes $-$ dans $-\frac{-6}{2\times(-1)}$ : on calcule d'abord
$2\times(-1)=-2$, puis $\frac{-6}{-2}=3$, et enfin $\alpha=-3$.}
\end{exo}

\begin{exo}[Le sommet à partir des racines]
\cas{1)}
\begin{calculs}
\Delta_{f} &= (-6)^{2}-4\times 1\times 8=4 &&\justif{$36-32$ ; $\sqrt{4}=2$}\\
x_{1} &= \frac{6-2}{2}=2 &&\\
x_{2} &= \frac{6+2}{2}=4 &&
\end{calculs}
\explique{Racines de $f$ : \res{$2$ et $4$}.}{}
\begin{calculs}
\Delta_{g} &= 8^{2}-4\times(-2)\times 10=144 &&\justif{$64+80$ ; $\sqrt{144}=12$}\\
x_{1} &= \frac{-8-12}{-4}=5 &&\justif{$2a=-4$}\\
x_{2} &= \frac{-8+12}{-4}=-1 &&
\end{calculs}
\explique{Racines de $g$ : \res{$-1$ et $5$}.}{}
\cas{2)}
\explique{Le sommet est sur l'axe de symétrie, au milieu des deux racines :
$\alpha=\dfrac{x_{1}+x_{2}}{2}$.}{remarque de la partie VII}
\begin{calculs}
\alpha_{f} &= \frac{2+4}{2}=\res{3} &&\\
\alpha_{g} &= \frac{-1+5}{2}=\res{2} &&
\end{calculs}
\cas{3)}
\begin{calculs}
\beta_{f} &= f(3)=3^{2}-6\times 3+8 &&\\
\beta_{f} &= 9-18+8=\res{-1} &&\\[3pt]
\beta_{g} &= g(2)=-2\times 2^{2}+8\times 2+10 &&\\
\beta_{g} &= -8+16+10=\res{18} &&
\end{calculs}
\explique{$f$ : $a=1>0$, minimum $-1$ en $x=3$.}{un creux}
\tabvar{f(x)}{$-\infty$,$3$,$+\infty$}{+/ ,-/$-1$,+/ }
\explique{$g$ : $a=-2<0$, maximum $18$ en $x=2$.}{une bosse}
\tabvar{g(x)}{$-\infty$,$2$,$+\infty$}{-/ ,+/$18$,-/ }
\rem{Contrôle avec la formule : $\alpha_{f}=-\frac{-6}{2}=3$ et
$\alpha_{g}=-\frac{8}{-4}=2$.}
\rem{Pour $g$, on peut aussi diviser $-2x^{2}+8x+10=0$ par $-2$ : $x^{2}-4x-5=0$ a les
mêmes racines.}
\end{exo}

\begin{exo}[Maximum ou minimum ?]
\explique{Le signe de $a$ dit « maximum » ou « minimum » ; la valeur est $\beta$, atteinte
en $x=\alpha$.}{théorème 3 du cours}
\cas{a)}
\begin{calculs}
\alpha &= -\frac{-8}{2\times 4}=1 &&\justif{$a=4$ et $b=-8$}\\
\beta &= 4\times 1^{2}-8\times 1+1=-3 &&
\end{calculs}
\explique{$a=4>0$ : $f_{1}$ a un \res{minimum, $-3$}, atteint en $x=1$.}{}
\cas{b)}
\begin{calculs}
\alpha &= -\frac{4}{2\times(-1)}=2 &&\justif{$a=-1$ et $b=4$}\\
\beta &= -2^{2}+4\times 2+2=6 &&
\end{calculs}
\explique{$a=-1<0$ : $f_{2}$ a un \res{maximum, $6$}, atteint en $x=2$.}{}
\cas{c)}
\begin{calculs}
\alpha &= -\frac{6}{2\times(-3)}=1 &&\justif{$a=-3$ et $b=6$}\\
\beta &= -3\times 1^{2}+6\times 1=3 &&
\end{calculs}
\explique{$a=-3<0$ : $f_{3}$ a un \res{maximum, $3$}, atteint en $x=1$.}{}
\cas{d)}
\begin{calculs}
\alpha &= -\frac{8}{2\times 1}=-4 &&\justif{$a=1$ et $b=8$}\\
\beta &= (-4)^{2}+8\times(-4)+12=-4 &&\justif{$16-32+12$}
\end{calculs}
\explique{$a=1>0$ : $f_{4}$ a un \res{minimum, $-4$}, atteint en $x=-4$.}{}
\explique{Ont un maximum : \res{$f_{2}$ et $f_{3}$} ; ont un minimum : \res{$f_{1}$ et
$f_{4}$}.}{on conclut par une phrase}
\rem{Le signe de $a$ suffit pour répondre « maximum » ou « minimum » ; le calcul de
$\beta$ donne sa valeur.}
\end{exo}

\partie{Synthèse : étudier et tracer une parabole}

\begin{exo}[Étude complète]
\explique{$a=1>0$ : la parabole est \res{un creux}.}{\textbf{sens}}
\begin{calculs}
\alpha &= -\frac{-2}{2\times 1}=1 &&\justif{\textbf{sommet}}\\
\beta &= 1^{2}-2\times 1-3=-4 &&\justif{$\beta=f(1)$ : $S(1\,;-4)$}\\
\Delta &= (-2)^{2}-4\times 1\times(-3)=16 &&\justif{\textbf{racines} ; $\sqrt{16}=4$}\\
x_{1} &= \frac{2-4}{2}=-1 &&\\
x_{2} &= \frac{2+4}{2}=3 &&
\end{calculs}
\explique{\textbf{Variations} sur $\intff{-2}{4}$ : $f(-2)=4+4-3=5$ et
$f(4)=16-8-3=5$.}{on calcule les images aux bornes}
\tabvar{f(x)}{$-2$,$1$,$4$}{+/$5$,-/$-4$,+/$5$}
\explique{$f$ a un \res{minimum, $-4$}, atteint en $x=1$.}{}
\begin{minipage}[c]{0.58\linewidth}\raggedright
\explique{\textbf{Valeurs}, avec $f(0)=-3$ et $f(2)=4-4-3=-3$ :}{symétrie par rapport à
$x=1$}
\begin{center}\small
\renewcommand{\arraystretch}{1.25}
\setlength{\tabcolsep}{4pt}
\begin{tabular}{|c|c|c|c|c|c|c|c|}
\hline
$x$ & $-2$ & $-1$ & $0$ & $1$ & $2$ & $3$ & $4$\\
\hline
$f(x)$ & $5$ & $0$ & $-3$ & $-4$ & $-3$ & $0$ & $5$\\
\hline
\end{tabular}
\end{center}
\explique{\textbf{Tracé} : on place les points et le sommet $S$, puis on relie par une
courbe arrondie, symétrique par rapport à la droite $x=1$.}{méthode 8 du cours}
\end{minipage}\hfill
\begin{minipage}[c]{0.38\linewidth}\centering
\begin{tikzpicture}[x=0.42cm,y=0.42cm]
  \repere{-3}{-5}{5}{6}
  \gradx{-2,-1,1,2,3,4} \grady{-4,-3,-2,-1,1,2,3,4,5}
  \draw[black!40,dashed] (1,-5) -- (1,6);
  \draw[coulCourbe,line width=1.1pt,domain=-2:4,samples=60] plot (\x,{\x*\x-2*\x-3});
  \foreach \p in {(-2,5),(-1,0),(0,-3),(2,-3),(3,0),(4,5)} \path[coulCourbe] \p node[croix]{};
  \path[coulExo] (1,-4) node[croix]{} node[below right,font=\scriptsize] {$S$};
  \node[coulCourbe,font=\footnotesize,right] at (4,5) {$\Cf$};
\end{tikzpicture}
\end{minipage}
\end{exo}

\begin{exo}[Étude complète]
\explique{$a=-1<0$ : la parabole est \res{une bosse}.}{\textbf{sens}}
\begin{calculs}
\alpha &= -\frac{-1}{2\times(-1)}=-\frac{1}{2} &&\justif{\textbf{sommet} ; $\frac{-1}{-2}=\frac{1}{2}$}\\
\beta &= -\left(-\frac{1}{2}\right)^{2}-\left(-\frac{1}{2}\right)+2 &&\justif{$\beta=g\left(-\frac{1}{2}\right)$}\\
\beta &= -\frac{1}{4}+\frac{2}{4}+\frac{8}{4}=\frac{9}{4} &&\justif{soit $2{,}25$ : $S\left(-\frac{1}{2}\,;\frac{9}{4}\right)$}\\
\Delta &= (-1)^{2}-4\times(-1)\times 2=9 &&\justif{\textbf{racines} ; $\sqrt{9}=3$}\\
x_{1} &= \frac{1-3}{-2}=1 &&\justif{$-b=1$ et $2a=-2$}\\
x_{2} &= \frac{1+3}{-2}=-2 &&
\end{calculs}
\explique{\textbf{Variations} sur $\intff{-3}{2}$ : $g(-3)=-9+3+2=-4$ et
$g(2)=-4-2+2=-4$.}{images aux bornes}
\tabvar{g(x)}{$-3$,$-\frac{1}{2}$,$2$}{-/$-4$,+/$\frac{9}{4}$,-/$-4$}
\explique{$g$ a un \res{maximum, $\dfrac{9}{4}$}, atteint en $x=-\dfrac{1}{2}$.}{}
\begin{minipage}[c]{0.58\linewidth}\raggedright
\explique{\textbf{Valeurs}, avec $g(-1)=-1+1+2=2$ et $g(0)=2$ :}{symétrie par rapport à
$x=-\frac{1}{2}$}
\begin{center}\small
\renewcommand{\arraystretch}{1.25}
\setlength{\tabcolsep}{4pt}
\begin{tabular}{|c|c|c|c|c|c|c|}
\hline
$x$ & $-3$ & $-2$ & $-1$ & $0$ & $1$ & $2$\\
\hline
$g(x)$ & $-4$ & $0$ & $2$ & $2$ & $0$ & $-4$\\
\hline
\end{tabular}
\end{center}
\explique{\textbf{Tracé} : le sommet est entre $x=-1$ et $x=0$, un peu au-dessus de
$2$.}{$\frac{9}{4}=2{,}25$}
\end{minipage}\hfill
\begin{minipage}[c]{0.38\linewidth}\centering
\begin{tikzpicture}[x=0.42cm,y=0.42cm]
  \repere{-4}{-5}{3}{4}
  \gradx{-3,-2,-1,1,2} \grady{-4,-3,-2,-1,1,2,3}
  \draw[black!40,dashed] (-0.5,-5) -- (-0.5,4);
  \draw[coulCourbe,line width=1.1pt,domain=-3:2,samples=60] plot (\x,{-\x*\x-\x+2});
  \foreach \p in {(-3,-4),(-2,0),(-1,2),(0,2),(1,0),(2,-4)} \path[coulCourbe] \p node[croix]{};
  \path[coulExo] (-0.5,2.25) node[croix]{} node[above,font=\scriptsize] {$S$};
  \node[coulCourbe,font=\footnotesize,right] at (2,-4) {$\Cg$};
\end{tikzpicture}
\end{minipage}
\end{exo}

\begin{exo}[Étude complète]
\explique{$a=2>0$ : la parabole est \res{un creux}.}{\textbf{sens}}
\begin{calculs}
\alpha &= -\frac{2}{2\times 2}=-\frac{1}{2} &&\justif{\textbf{sommet}}\\
\beta &= 2\times\left(-\frac{1}{2}\right)^{2}+2\times\left(-\frac{1}{2}\right)-4 &&\justif{$\beta=h\left(-\frac{1}{2}\right)$}\\
\beta &= \frac{1}{2}-1-4=-\frac{9}{2} &&\justif{$2\times\frac{1}{4}=\frac{1}{2}$ ; soit $-4{,}5$}\\
\Delta &= 2^{2}-4\times 2\times(-4)=36 &&\justif{\textbf{racines} ; $\sqrt{36}=6$}\\
x_{1} &= \frac{-2-6}{4}=-2 &&\\
x_{2} &= \frac{-2+6}{4}=1 &&
\end{calculs}
\explique{\textbf{Variations} sur $\intff{-2}{1}$ : les bornes sont les racines, donc
$h(-2)=h(1)=0$.}{}
\tabvar{h(x)}{$-2$,$-\frac{1}{2}$,$1$}{+/$0$,-/$-\frac{9}{2}$,+/$0$}
\explique{$h$ a un \res{minimum, $-\dfrac{9}{2}$}, atteint en $x=-\dfrac{1}{2}$.}{}
\begin{minipage}[c]{0.58\linewidth}\raggedright
\explique{\textbf{Valeurs}, avec $h(-1)=2-2-4=-4$ et $h(0)=-4$ :}{}
\begin{center}\small
\renewcommand{\arraystretch}{1.25}
\begin{tabular}{|c|c|c|c|c|}
\hline
$x$ & $-2$ & $-1$ & $0$ & $1$\\
\hline
$h(x)$ & $0$ & $-4$ & $-4$ & $0$\\
\hline
\end{tabular}
\end{center}
\explique{La courbe passe bien par les points $(-2\,;0)$ et $(1\,;0)$ de l'axe des
abscisses : ce sont les racines calculées.}{la vérification demandée}
\end{minipage}\hfill
\begin{minipage}[c]{0.38\linewidth}\centering
\begin{tikzpicture}[x=0.5cm,y=0.5cm]
  \repere{-3}{-5}{2}{2}
  \gradx{-2,-1,1} \grady{-4,-3,-2,-1,1}
  \draw[black!40,dashed] (-0.5,-5) -- (-0.5,2);
  \draw[coulCourbe,line width=1.1pt,domain=-2:1,samples=60] plot (\x,{2*\x*\x+2*\x-4});
  \foreach \p in {(-2,0),(-1,-4),(0,-4),(1,0)} \path[coulCourbe] \p node[croix]{};
  \path[coulExo] (-0.5,-4.5) node[croix]{} node[below,font=\scriptsize] {$S$};
  \node[coulCourbe,font=\footnotesize,above right] at (1,0) {$\mathcal{C}_{h}$};
\end{tikzpicture}
\end{minipage}
\end{exo}

\partie{Signe d'un trinôme}

\begin{exo}[À partir de la forme factorisée]
\explique{On dresse le tableau de signes du produit : une ligne par facteur, rangée dans
l'ordre des racines, puis la règle des signes dans chaque colonne.}{comme l'exemple de la
partie IX}
\cas{$f(x)=3(x+2)(x-5)$}
\explique{$3$ est toujours positif ; $x+2$ s'annule en $-2$ et $x-5$ en $5$, négatifs
avant, positifs après.}{un facteur $x-r$ est négatif avant $r$, positif après}
\begin{center}
\begin{tikzpicture}[scale=0.85]
  \tkzTabInit[lgt=2.4,espcl=2.2]{$x$/1,$3$/1,$x+2$/1,$x-5$/1,$f(x)$/1.2}%
    {$-\infty$,$-2$,$5$,$+\infty$}
  \tkzTabLine{,+,t,+,t,+,}
  \tkzTabLine{,-,z,+,t,+,}
  \tkzTabLine{,-,t,-,z,+,}
  \tkzTabLine{,+,z,-,z,+,}
\end{tikzpicture}
\end{center}
\cas{$g(x)=-2(x-1)(x+6)$}
\explique{On range les racines : $-6<1$. Le facteur $-2$ est toujours négatif.}{la ligne
de $x+6$ vient en premier}
\begin{center}
\begin{tikzpicture}[scale=0.85]
  \tkzTabInit[lgt=2.4,espcl=2.2]{$x$/1,$-2$/1,$x+6$/1,$x-1$/1,$g(x)$/1.2}%
    {$-\infty$,$-6$,$1$,$+\infty$}
  \tkzTabLine{,-,t,-,t,-,}
  \tkzTabLine{,-,z,+,t,+,}
  \tkzTabLine{,-,t,-,z,+,}
  \tkzTabLine{,-,z,+,z,-,}
\end{tikzpicture}
\end{center}
\cas{$h(x)=4(x+4)(3x-1)$}
\explique{$3x-1=0$ donne $x=\dfrac{1}{3}$ ; $3x-1$ est négatif avant $\dfrac{1}{3}$,
positif après.}{coefficient de $x$ positif : $3>0$}
\begin{center}
\begin{tikzpicture}[scale=0.85]
  \tkzTabInit[lgt=2.4,espcl=2.2]{$x$/1,$4$/1,$x+4$/1,$3x-1$/1,$h(x)$/1.2}%
    {$-\infty$,$-4$,$\frac{1}{3}$,$+\infty$}
  \tkzTabLine{,+,t,+,t,+,}
  \tkzTabLine{,-,z,+,t,+,}
  \tkzTabLine{,-,t,-,z,+,}
  \tkzTabLine{,+,z,-,z,+,}
\end{tikzpicture}
\end{center}
\cas{$k(x)=-3(2x-5)(4x+1)$}
\explique{$2x-5=0$ donne $x=\dfrac{5}{2}$ ; $4x+1=0$ donne $x=-\dfrac{1}{4}$. On range :
$-\dfrac{1}{4}<\dfrac{5}{2}$.}{la ligne de $4x+1$ vient en premier}
\begin{center}
\begin{tikzpicture}[scale=0.85]
  \tkzTabInit[lgt=2.4,espcl=2.2]{$x$/1,$-3$/1,$4x+1$/1,$2x-5$/1,$k(x)$/1.2}%
    {$-\infty$,$-\frac{1}{4}$,$\frac{5}{2}$,$+\infty$}
  \tkzTabLine{,-,t,-,t,-,}
  \tkzTabLine{,-,z,+,t,+,}
  \tkzTabLine{,-,t,-,z,+,}
  \tkzTabLine{,-,z,+,z,-,}
\end{tikzpicture}
\end{center}
\rem{Contrôle par le théorème 4 : le coefficient de $x^{2}$ est $a=3$ pour $f$, $-2$ pour
$g$, $4\times 3=12$ pour $h$ et $-3\times 2\times 4=-24$ pour $k$. Le trinôme a le signe
de $a$ à l'extérieur des racines : c'est bien la dernière ligne de chaque tableau.}
\end{exo}

\begin{exo}[À partir de la forme développée]
\explique{On calcule $\Delta$ et les racines ; le trinôme a le signe de $a$ à l'extérieur
des racines, le signe contraire entre elles.}{méthode 9 et théorème 4 du cours}
\cas{a)}
\begin{calculs}
\Delta &= 5^{2}-4\times 2\times(-12)=121 &&\justif{$25+96$ ; $\sqrt{121}=11$}\\
x_{1} &= \frac{-5-11}{4}=-4 &&\\
x_{2} &= \frac{-5+11}{4}=\frac{3}{2} &&
\end{calculs}
\tabsignes{f(x)}{$-\infty$,$-4$,$\frac{3}{2}$,$+\infty$}{,+,z,-,z,+,}
\explique{$a=2>0$ : positif à l'extérieur des racines, négatif entre.}{}
\cas{b)}
\begin{calculs}
\Delta &= 2^{2}-4\times(-1)\times 8=36 &&\justif{$4+32$ ; $\sqrt{36}=6$}\\
x_{1} &= \frac{-2-6}{-2}=4 &&\\
x_{2} &= \frac{-2+6}{-2}=-2 &&\justif{on range les racines : $-2<4$}
\end{calculs}
\tabsignes{g(x)}{$-\infty$,$-2$,$4$,$+\infty$}{,-,z,+,z,-,}
\explique{$a=-1<0$ : négatif à l'extérieur, positif entre les racines.}{}
\cas{c)}
\begin{calculs}
\Delta &= (-7)^{2}-4\times 2\times 3=25 &&\justif{$49-24$ ; $\sqrt{25}=5$}\\
x_{1} &= \frac{7-5}{4}=\frac{1}{2} &&\\
x_{2} &= \frac{7+5}{4}=3 &&
\end{calculs}
\tabsignes{h(x)}{$-\infty$,$\frac{1}{2}$,$3$,$+\infty$}{,+,z,-,z,+,}
\explique{$a=2>0$ : positif à l'extérieur des racines.}{}
\cas{d)}
\begin{calculs}
\Delta &= (-1)^{2}-4\times(-3)\times 2=25 &&\justif{$1+24$ ; $\sqrt{25}=5$}\\
x_{1} &= \frac{1-5}{-6}=\frac{2}{3} &&\justif{$-b=1$ et $2a=-6$}\\
x_{2} &= \frac{1+5}{-6}=-1 &&\justif{on range : $-1<\frac{2}{3}$}
\end{calculs}
\tabsignes{k(x)}{$-\infty$,$-1$,$\frac{2}{3}$,$+\infty$}{,-,z,+,z,-,}
\explique{$a=-3<0$ : négatif à l'extérieur, positif entre les racines.}{}
\rem{Avec $a<0$, la formule donne souvent la plus grande racine en premier : on range
toujours les racines dans l'ordre croissant avant de remplir le tableau.}
\end{exo}

\begin{exo}[Les trois cas]
\explique{$\Delta>0$ : signe de $a$ à l'extérieur des racines ; $\Delta=0$ : signe de
$a$, nul en $x_{0}$ ; $\Delta<0$ : toujours le signe de $a$.}{théorème 4 du cours}
\cas{$f(x)=x^{2}+4x+4$}
\begin{calculs}
\Delta &= 4^{2}-4\times 1\times 4=0 &&\justif{une racine double}\\
x_{0} &= \frac{-4}{2}=-2 &&
\end{calculs}
\tabsignes{f(x)}{$-\infty$,$-2$,$+\infty$}{,+,z,+,}
\explique{$a=1>0$ : $f(x)$ est \res{positif, nul en $-2$}.}{$f(x)=(x+2)^{2}$, un carré}
\cas{$g(x)=x^{2}-4$}
\explique{$x^{2}-4=(x-2)(x+2)$ : racines $-2$ et $2$ ; $a=1>0$.}{différence de carrés ; ou
$\Delta=16$}
\tabsignes{g(x)}{$-\infty$,$-2$,$2$,$+\infty$}{,+,z,-,z,+,}
\explique{$g(x)$ est \res{positif à l'extérieur de $-2$ et $2$, négatif entre}.}{}
\cas{$h(x)=x^{2}+4$}
\begin{calculs}
\Delta &= 0^{2}-4\times 1\times 4=-16 &&\justif{$\Delta<0$ : pas de racine}
\end{calculs}
\tabsignes[4]{h(x)}{$-\infty$,$+\infty$}{,+,}
\explique{$a=1>0$ : $h(x)$ est \res{toujours positif}.}{}
\cas{$k(x)=-x^{2}+2x-4$}
\begin{calculs}
\Delta &= 2^{2}-4\times(-1)\times(-4)=-12 &&\justif{$4-16$ : pas de racine}
\end{calculs}
\tabsignes[4]{k(x)}{$-\infty$,$+\infty$}{,-,}
\explique{$a=-1<0$ : $k(x)$ est \res{toujours négatif}.}{}
\cas{$l(x)=-9x^{2}+6x-1$}
\begin{calculs}
\Delta &= 6^{2}-4\times(-9)\times(-1)=0 &&\justif{$36-36$}\\
x_{0} &= \frac{-6}{-18}=\frac{1}{3} &&
\end{calculs}
\tabsignes{l(x)}{$-\infty$,$\frac{1}{3}$,$+\infty$}{,-,z,-,}
\explique{$a=-9<0$ : $l(x)$ est \res{négatif, nul en $\dfrac{1}{3}$}.}{}
\cas{$m(x)=2x^{2}-5x+3$}
\begin{calculs}
\Delta &= (-5)^{2}-4\times 2\times 3=1 &&\justif{$25-24$}\\
x_{1} &= \frac{5-1}{4}=1 &&\\
x_{2} &= \frac{5+1}{4}=\frac{3}{2} &&
\end{calculs}
\tabsignes{m(x)}{$-\infty$,$1$,$\frac{3}{2}$,$+\infty$}{,+,z,-,z,+,}
\explique{$a=2>0$ : $m(x)$ est \res{positif à l'extérieur de $1$ et $\dfrac{3}{2}$, négatif
entre}.}{}
\rem{Erreur fréquente : quand $\Delta<0$, conclure « toujours positif » sans regarder
$a$. Pour $k$, $a<0$ : toujours négatif.}
\end{exo}

\begin{exo}[Sans tableau de signes]
\explique{Si $\Delta<0$, le trinôme a toujours le signe de $a$ ; si $\Delta>0$, il change
de signe en chaque racine ; si $\Delta=0$, il a le signe de $a$ et s'annule une
fois.}{théorème 4 du cours}
\begin{calculs}
\Delta_{f} &= 2^{2}-4\times 1\times 3=-8 &&\justif{$\Delta<0$ et $a=1>0$ : \res{toujours positif}}\\
\Delta_{g} &= 0^{2}-4\times(-2)\times(-3)=-24 &&\justif{$\Delta<0$ et $a=-2<0$ : \res{toujours négatif}}\\
\Delta_{h} &= 0^{2}-4\times 3\times(-12)=144 &&\justif{$\Delta>0$ : \res{change de signe} (en $-2$ et $2$)}\\
\Delta_{k} &= 4^{2}-4\times(-1)\times(-4)=0 &&\justif{$\Delta=0$ et $a=-1<0$ : \res{négatif}, nul en $2$}\\
\Delta_{l} &= 1^{2}-4\times 4\times 1=-15 &&\justif{$\Delta<0$ et $a=4>0$ : \res{toujours positif}}\\
\Delta_{m} &= 2^{2}-4\times(-3)\times(-1)=-8 &&\justif{$\Delta<0$ et $a=-3<0$ : \res{toujours négatif}}
\end{calculs}
\rem{$k(x)=-(x-2)^{2}$ ne change pas de signe : il est négatif, sauf en $x=2$ où il est
nul. On peut dire « toujours négatif ou nul ».}
\end{exo}

\partie{Inéquations du second degré}

\begin{exo}[Inéquations]
\explique{On cherche les racines, on dresse le tableau de signes, puis on lit les $x$ qui
conviennent : crochets ouverts pour $<$ ou $>$, fermés aux racines pour $\le$ ou
$\ge$.}{méthode 10 du cours}
\cas{a)}
\begin{calculs}
\Delta &= 4^{2}-4\times 1\times 4=0 &&\justif{une racine double}\\
x_{0} &= \frac{-4}{2}=-2 &&
\end{calculs}
\tabsignes{x^{2}+4x+4}{$-\infty$,$-2$,$+\infty$}{,+,z,+,}
\explique{$\res{S=\left]-\infty\,;-2\right[\cup\left]-2\,;+\infty\right[}$.}{strictement
positif partout, sauf en $-2$ où il est nul}
\cas{b)}
\begin{calculs}
\Delta &= 5^{2}-4\times(-2)\times(-3)=1 &&\justif{$25-24$}\\
x_{1} &= \frac{-5-1}{-4}=\frac{3}{2} &&\\
x_{2} &= \frac{-5+1}{-4}=1 &&\justif{on range : $1<\frac{3}{2}$}
\end{calculs}
\tabsignes{-2x^{2}+5x-3}{$-\infty$,$1$,$\frac{3}{2}$,$+\infty$}{,-,z,+,z,-,}
\explique{$\res{S=\left]-\infty\,;1\right]\cup\left[\tfrac{3}{2}\,;+\infty\right[}$.}{$a<0$
: négatif à l'extérieur ; $\le$ : racines comprises}
\cas{c)}
\begin{calculs}
\Delta &= (-2)^{2}-4\times 1\times(-15)=64 &&\justif{$\sqrt{64}=8$}\\
x_{1} &= \frac{2-8}{2}=-3 &&\\
x_{2} &= \frac{2+8}{2}=5 &&
\end{calculs}
\tabsignes{x^{2}-2x-15}{$-\infty$,$-3$,$5$,$+\infty$}{,+,z,-,z,+,}
\explique{$\res{S=\left]-\infty\,;-3\right]\cup\left[5\,;+\infty\right[}$.}{$a>0$ :
positif à l'extérieur ; $\ge$ : racines comprises}
\cas{d)}
\begin{calculs}
\Delta &= 3^{2}-4\times 1\times(-10)=49 &&\justif{$9+40$ ; $\sqrt{49}=7$}\\
x_{1} &= \frac{-3-7}{2}=-5 &&\\
x_{2} &= \frac{-3+7}{2}=2 &&
\end{calculs}
\tabsignes{x^{2}+3x-10}{$-\infty$,$-5$,$2$,$+\infty$}{,+,z,-,z,+,}
\explique{$\res{S=\left]-5\,;2\right[}$.}{négatif entre les racines ; $<$ : racines
exclues}
\rem{En a), on peut aussi écrire $S=\R\setminus\{-2\}$ (« $\R$ privé de $-2$ »).}
\end{exo}

\begin{exo}[Tout ramener d'un côté]
\explique{On met tout dans le membre de gauche, puis on résout comme à l'exercice
36.}{méthode 10 du cours, cas c)}
\cas{a)}
\begin{calculs}
-2x^{2}+7x &\le 3 &&\justif{on recopie}\\
-2x^{2}+7x-3 &\le 0 &&\justif{on retranche $3$}
\end{calculs}
\begin{calculs}
\Delta &= 7^{2}-4\times(-2)\times(-3)=25 &&\justif{$49-24$ ; $\sqrt{25}=5$}\\
x_{1} &= \frac{-7-5}{-4}=3 &&\\
x_{2} &= \frac{-7+5}{-4}=\frac{1}{2} &&
\end{calculs}
\tabsignes{-2x^{2}+7x-3}{$-\infty$,$\frac{1}{2}$,$3$,$+\infty$}{,-,z,+,z,-,}
\explique{$\res{S=\left]-\infty\,;\tfrac{1}{2}\right]\cup\left[3\,;+\infty\right[}$.}{$a<0$
: négatif à l'extérieur}
\cas{b)}
\begin{calculs}
2x^{2} &\ge 3x-4 &&\justif{on recopie}\\
2x^{2}-3x+4 &\ge 0 &&\justif{on retranche $3x$ et on ajoute $4$}
\end{calculs}
\begin{calculs}
\Delta &= (-3)^{2}-4\times 2\times 4=-23 &&\justif{$9-32$ : pas de racine}
\end{calculs}
\explique{$a=2>0$ : le trinôme est toujours positif ; $\res{S=\R}$.}{signe de $a$
partout}
\cas{c)}
\begin{calculs}
x(x-4) &\ge x-6 &&\justif{on recopie}\\
x^{2}-4x &\ge x-6 &&\justif{on développe}\\
x^{2}-5x+6 &\ge 0 &&\justif{on retranche $x$ et on ajoute $6$}
\end{calculs}
\begin{calculs}
\Delta &= (-5)^{2}-4\times 1\times 6=1 &&\\
x_{1} &= \frac{5-1}{2}=2 &&\\
x_{2} &= \frac{5+1}{2}=3 &&
\end{calculs}
\tabsignes{x^{2}-5x+6}{$-\infty$,$2$,$3$,$+\infty$}{,+,z,-,z,+,}
\explique{$\res{S=\left]-\infty\,;2\right]\cup\left[3\,;+\infty\right[}$.}{$a>0$ :
positif à l'extérieur}
\cas{d)}
\begin{calculs}
x^{2}+8 &< 6x &&\justif{on recopie}\\
x^{2}-6x+8 &< 0 &&\justif{on retranche $6x$}
\end{calculs}
\begin{calculs}
\Delta &= (-6)^{2}-4\times 1\times 8=4 &&\justif{$36-32$ ; $\sqrt{4}=2$}\\
x_{1} &= \frac{6-2}{2}=2 &&\\
x_{2} &= \frac{6+2}{2}=4 &&
\end{calculs}
\tabsignes{x^{2}-6x+8}{$-\infty$,$2$,$4$,$+\infty$}{,+,z,-,z,+,}
\explique{$\res{S=\left]2\,;4\right[}$.}{négatif entre les racines ; inégalité stricte}
\rem{On ne divise jamais les deux membres par $x$ : on ne connaît pas son signe.}
\end{exo}

\begin{exo}[Cas particuliers]
\cas{a)}
\begin{calculs}
x^{2} &\ge 25 &&\justif{on recopie}\\
x^{2}-25 &\ge 0 &&\justif{tout dans le membre de gauche}\\
(x-5)(x+5) &\ge 0 &&\justif{différence de carrés : racines $-5$ et $5$}
\end{calculs}
\explique{$a=1>0$ : positif à l'extérieur des racines ;
$\res{S=\left]-\infty\,;-5\right]\cup\left[5\,;+\infty\right[}$.}{}
\cas{b)}
\begin{calculs}
x^{2} &< 9 &&\justif{on recopie}\\
(x-3)(x+3) &< 0 &&\justif{$x^{2}-9<0$ : racines $-3$ et $3$}
\end{calculs}
\explique{Négatif entre les racines : $\res{S=\left]-3\,;3\right[}$.}{$a=1>0$}
\cas{c)}
\explique{$-x^{2}+1=(1-x)(1+x)$ : racines $-1$ et $1$ ; $a=-1<0$, donc positif entre
les racines.}{}
\explique{$\res{S=\left]-1\,;1\right[}$.}{inégalité stricte}
\cas{d)}
\explique{$3x^{2}\ge 0$, donc $3x^{2}+6\ge 6$ : le membre de gauche n'est jamais
négatif ni nul ; $\res{S=\varnothing}$.}{ou : $\Delta=-72<0$ et $a=3>0$, toujours positif}
\cas{e)}
\explique{Un carré est positif ; $(x+2)^{2}$ n'est nul qu'en $-2$.
$\res{S=\left]-\infty\,;-2\right[\cup\left]-2\,;+\infty\right[}$.}{inégalité stricte :
on enlève $-2$}
\cas{f)}
\begin{calculs}
\Delta &= (-1)^{2}-4\times 1\times 4=-15 &&\justif{$\Delta<0$ : pas de racine}
\end{calculs}
\explique{$a=1>0$ : toujours positif ; $\res{S=\R}$.}{}
\rem{Erreur fréquente en a) : écrire $x\ge 5$ seulement. Les nombres négatifs comme
$-6$ conviennent aussi : $(-6)^{2}=36\ge 25$.}
\end{exo}

\begin{exo}[Lecture graphique]
\cas{a)}
\begin{minipage}[c]{0.6\linewidth}\raggedright
\explique{Graphiquement : la courbe coupe l'axe des abscisses en $-3$ et en $1$.}{les
points où $f(x)=0$}
\begin{calculs}
\Delta &= 2^{2}-4\times 1\times(-3)=16 &&\justif{$\sqrt{16}=4$}\\
x_{1} &= \frac{-2-4}{2}=-3 &&\\
x_{2} &= \frac{-2+4}{2}=1 &&
\end{calculs}
\explique{$\res{S=\ens{-3\,;1}}$ : le calcul confirme la lecture.}{}
\end{minipage}\hfill
\begin{minipage}[c]{0.36\linewidth}\centering
\begin{tikzpicture}[x=0.32cm,y=0.32cm]
  \repere{-5}{-5}{3}{5}
  \gradx{-3,1} \grady{-4,4}
  \draw[coulCourbe,line width=1.1pt,domain=-4:2,samples=60] plot (\x,{\x*\x+2*\x-3});
  \path[coulInt] (-3,0) node[croix]{} (1,0) node[croix]{};
\end{tikzpicture}
\end{minipage}
\cas{b)}
\begin{minipage}[c]{0.6\linewidth}\raggedright
\explique{Graphiquement : la courbe est au-dessus de l'axe des abscisses à gauche de $-3$
et à droite de $1$.}{$f(x)>0$ : points au-dessus de l'axe}
\explique{Par le calcul : $a=1>0$, donc $f(x)$ est positif à l'extérieur des racines.}{théorème 4}
\explique{$\res{S=\left]-\infty\,;-3\right[\cup\left]1\,;+\infty\right[}$.}{racines exclues}
\end{minipage}\hfill
\begin{minipage}[c]{0.36\linewidth}\centering
\begin{tikzpicture}[x=0.32cm,y=0.32cm]
  \repere{-5}{-5}{3}{5}
  \gradx{-3,1} \grady{-4,4}
  \draw[coulCourbe,line width=1.1pt,domain=-4:2,samples=60] plot (\x,{\x*\x+2*\x-3});
  \morceau{-5}{-3} \morceau[coulIntB]{1}{3}
\end{tikzpicture}
\end{minipage}
\cas{c)}
\begin{minipage}[c]{0.6\linewidth}\raggedright
\explique{Graphiquement : la courbe est sous l'axe ou sur l'axe, entre $-3$ et $1$.}{$f(x)\le 0$ :
« ou nul », racines comprises}
\explique{Par le calcul : négatif entre les racines ; $\res{S=\intff{-3}{1}}$.}{}
\end{minipage}\hfill
\begin{minipage}[c]{0.36\linewidth}\centering
\begin{tikzpicture}[x=0.32cm,y=0.32cm]
  \repere{-5}{-5}{3}{5}
  \gradx{-3,1} \grady{-4,4}
  \draw[coulCourbe,line width=1.1pt,domain=-4:2,samples=60] plot (\x,{\x*\x+2*\x-3});
  \morceau{-3}{1}
\end{tikzpicture}
\end{minipage}
\cas{d)}
\begin{minipage}[c]{0.6\linewidth}\raggedright
\explique{Graphiquement : on trace la droite horizontale $y=-3$ ; elle coupe la courbe en
deux points, d'abscisses $-2$ et $0$.}{les antécédents de $-3$}
\begin{calculs}
x^{2}+2x-3 &= -3 &&\justif{on recopie $f(x)=-3$}\\
x^{2}+2x &= 0 &&\justif{on ajoute $3$}\\
x(x+2) &= 0 &&\justif{facteur commun $x$}\\
x &= 0 \quad\text{ou}\quad x+2=0 &&\justif{produit nul}
\end{calculs}
\explique{$\res{S=\ens{-2\,;0}}$.}{}
\end{minipage}\hfill
\begin{minipage}[c]{0.36\linewidth}\centering
\begin{tikzpicture}[x=0.32cm,y=0.32cm]
  \repere{-5}{-5}{3}{5}
  \gradx{-2,1} \grady{-3,4}
  \draw[coulCourbe,line width=1.1pt,domain=-4:2,samples=60] plot (\x,{\x*\x+2*\x-3});
  \draw[coulCourbeB,line width=0.9pt] (-5,-3) -- (3,-3);
  \draw[coulInt,dashed] (-2,-3) -- (-2,0);
  \path[coulInt] (-2,-3) node[croix]{} (0,-3) node[croix]{};
\end{tikzpicture}
\end{minipage}
\rem{En b), la courbe ne montre que l'intervalle $\intff{-4}{2}$ ; c'est le calcul qui dit
que $f(x)$ reste positif au-delà.}
\end{exo}

\partie{Intersection d'une parabole et d'une droite}

\begin{exo}[Points d'intersection]
\explique{On égale les deux expressions de $y$ (parabole et droite) : les solutions sont les
abscisses des points communs. On met tout à gauche et on résout. Les ordonnées se calculent avec l'équation de la
droite.}{méthode 11 du cours}
\cas{a)}
\begin{calculs}
x^{2}-4x+2 &= -x+2 &&\justif{on égale les deux expressions de $y$}\\
x^{2}-3x &= 0 &&\justif{tout à gauche}\\
x(x-3) &= 0 &&\justif{pas de nombre seul : facteur commun $x$}\\
x &= 0 \quad\text{ou}\quad x=3 &&\justif{produit nul}
\end{calculs}
\explique{Ordonnées avec la droite : $0$ donne $2$ et $3$ donne $-3+2=-1$. Points \res{$(0\,;2)$ et
$(3\,;-1)$}.}{contrôle sur la parabole : $9-12+2=-1$}
\cas{b)}
\begin{calculs}
x^{2}+2x+4 &= 4x+3 &&\justif{on égale les deux expressions de $y$}\\
x^{2}-2x+1 &= 0 &&\justif{tout à gauche}\\
\Delta &= (-2)^{2}-4\times 1\times 1=0 &&\justif{une seule solution}\\
x_{0} &= \frac{2}{2}=1 &&
\end{calculs}
\explique{Ordonnée : $4\times 1+3=7$. Un seul point, \res{$(1\,;7)$} : la droite est
tangente à la parabole.}{elle la touche sans la traverser}
\cas{c)}
\begin{calculs}
x^{2}-x+3 &= x+1 &&\justif{on égale les deux expressions de $y$}\\
x^{2}-2x+2 &= 0 &&\justif{tout à gauche}\\
\Delta &= (-2)^{2}-4\times 1\times 2=-4 &&\justif{$\Delta<0$ : aucune solution}
\end{calculs}
\explique{\res{Aucun point commun} : la droite ne coupe pas la parabole.}{}
\rem{Erreur fréquente : donner seulement les abscisses. Un point a deux coordonnées.}
\end{exo}

\begin{exo}[Deux paraboles]
\explique{Même méthode : on égale les deux expressions de $y$ ; les solutions sont les
abscisses des points communs.}{propriété 5 du cours}
\cas{a)}
\begin{calculs}
x^{2}+x-3 &= -x^{2}+3x+1 &&\justif{on égale les deux expressions}\\
2x^{2}-2x-4 &= 0 &&\justif{tout à gauche : $x^{2}+x^{2}$ ; $x-3x$ ; $-3-1$}\\
\Delta &= (-2)^{2}-4\times 2\times(-4)=36 &&\justif{$4+32$ ; $\sqrt{36}=6$}\\
x_{1} &= \frac{2-6}{4}=-1 &&\\
x_{2} &= \frac{2+6}{4}=2 &&
\end{calculs}
\explique{Ordonnées avec $\mathcal{P}$ : $(-1)^{2}-1-3=-3$ et $2^{2}+2-3=3$. Points
\res{$(-1\,;-3)$ et $(2\,;3)$}.}{contrôle avec $\mathcal{P}'$ : $-1-3+1=-3$}
\cas{b)}
\begin{calculs}
x^{2}-2x+4 &= -x^{2}+2x+2 &&\\
2x^{2}-4x+2 &= 0 &&\justif{tout à gauche}\\
\Delta &= (-4)^{2}-4\times 2\times 2=0 &&\justif{une seule solution}\\
x_{0} &= \frac{4}{4}=1 &&
\end{calculs}
\explique{Ordonnée : $1-2+4=3$. Un seul point, \res{$(1\,;3)$} : les deux paraboles se
touchent.}{}
\cas{c)}
\begin{calculs}
3x^{2}-2x+1 &= 3x^{2}+4x-5 &&\\
-6x+6 &= 0 &&\justif{les $3x^{2}$ s'éliminent}\\
-6x &= -6 &&\justif{on retranche $6$}\\
x &= 1 &&\justif{on divise par $-6$}
\end{calculs}
\explique{Ordonnée : $3-2+1=2$. Un seul point, \res{$(1\,;2)$}.}{}
\explique{On remarque que les termes en $x^{2}$ s'éliminent : l'équation est du
\res{premier degré}, et il n'y a qu'un point commun.}{même coefficient $a=3$}
\end{exo}

\begin{exo}[Combien de points ?]
\explique{Pour chaque droite, on égale $x^{2}$ et l'expression de $y$ : les solutions
sont les abscisses des points communs ; le signe
du discriminant donne leur nombre, sans calculer les points.}{propriété 5 du cours}
\cas{$\mathcal{D}_{1}$}
\begin{calculs}
x^{2} &= 4x-4 &&\\
x^{2}-4x+4 &= 0 &&\justif{tout à gauche}\\
\Delta &= (-4)^{2}-4\times 1\times 4=0 &&
\end{calculs}
\explique{\res{Un seul point commun} : $\mathcal{D}_{1}$ est tangente à
$\mathcal{P}$.}{$\Delta=0$}
\cas{$\mathcal{D}_{2}$}
\begin{calculs}
x^{2} &= 4x-5 &&\\
x^{2}-4x+5 &= 0 &&\\
\Delta &= (-4)^{2}-4\times 1\times 5=-4 &&
\end{calculs}
\explique{\res{Aucun point commun}.}{$\Delta<0$}
\cas{$\mathcal{D}_{3}$}
\begin{calculs}
x^{2} &= 4x+1 &&\\
x^{2}-4x-1 &= 0 &&\\
\Delta &= (-4)^{2}-4\times 1\times(-1)=20 &&
\end{calculs}
\explique{\res{Deux points communs}.}{$\Delta>0$}
\rem{Les trois droites sont parallèles (même coefficient directeur $4$) : la plus basse
ne touche pas la parabole, la plus haute la coupe deux fois.}
\end{exo}

\begin{exo}[Avec un paramètre]
\cas{1)}
\begin{calculs}
x^{2}+x &= 3x+m &&\justif{même ordonnée pour la même abscisse}\\
x^{2}+x-3x-m &= 0 &&\justif{tout à gauche}\\
x^{2}-2x-m &= 0 &&
\end{calculs}
\explique{Les abscisses des points communs sont bien les solutions de
$\res{x^{2}-2x-m=0}$.}{}
\cas{2)}
\begin{calculs}
\Delta &= (-2)^{2}-4\times 1\times(-m) &&\justif{$a=1$, $b=-2$, $c=-m$}\\
\Delta &= \res{4+4m} &&\justif{$-4\times(-m)=+4m$}
\end{calculs}
\cas{3)}
\explique{Deux points quand $\Delta>0$ :}{}
\begin{calculs}
4+4m &> 0 &&\\
4m &> -4 &&\justif{on retranche $4$}\\
m &> -1 &&\justif{on divise par $4>0$ : le sens ne change pas}
\end{calculs}
\explique{La droite coupe $\mathcal{P}$ en \res{deux points si $m>-1$}, en \res{un seul si
$m=-1$} et en \res{aucun si $m<-1$}.}{$\Delta=0$ pour $m=-1$ ; $\Delta<0$ pour $m<-1$}
\rem{Erreur fréquente : prendre $c=m$. Après avoir tout mis à gauche, le nombre seul est
$-m$.}
\end{exo}

\partie{Position relative de deux courbes}

\begin{exo}[Une parabole et une droite]
\cas{1)}
\begin{calculs}
d(x) &= x^{2}+x-2-(3x-2) &&\justif{on soustrait \textbf{toute} l'expression de $g$}\\
d(x) &= x^{2}+x-2-3x+2 &&\justif{le $-$ change tous les signes}\\
d(x) &= \res{x^{2}-2x} &&
\end{calculs}
\cas{2)}
\begin{calculs}
d(x) &= x(x-2) &&\justif{facteur commun $x$ : racines $0$ et $2$}
\end{calculs}
\tabsignes{d(x)}{$-\infty$,$0$,$2$,$+\infty$}{,+,z,-,z,+,}
\explique{$a=1>0$ : positif à l'extérieur des racines, négatif entre.}{théorème 4}
\cas{3)}
\begin{minipage}[c]{0.6\linewidth}\raggedright
\explique{$\Cf$ est \res{au-dessus} de $\mathcal{C}_{g}$ sur $\left]-\infty\,;0\right[$
et sur $\left]2\,;+\infty\right[$, \res{au-dessous} sur $\left]0\,;2\right[$.}{là où
$d(x)>0$, puis là où $d(x)<0$}
\explique{Elles se coupent en $(0\,;-2)$ et en $(2\,;4)$.}{ordonnées : $g(0)=-2$ et
$g(2)=4$}
\end{minipage}\hfill
\begin{minipage}[c]{0.36\linewidth}\centering
\begin{tikzpicture}[x=0.42cm,y=0.42cm]
  \repere{-3}{-3}{3}{5}
  \gradx{-2,-1,1,2} \grady{-2,-1,1,2,3,4}
  \draw[coulCourbe,line width=1.1pt,domain=-3:2.3,samples=60] plot (\x,{\x*\x+\x-2});
  \draw[coulCourbeB,line width=1.1pt,domain=-0.33:2.33] plot (\x,{3*\x-2});
  \path (0,-2) node[croix]{} (2,4) node[croix]{};
  \node[coulCourbe,font=\footnotesize,left] at (-2.6,3) {$\Cf$};
  \node[coulCourbeB,font=\footnotesize,right] at (2.3,4.6) {$\mathcal{C}_{g}$};
\end{tikzpicture}
\end{minipage}
\end{exo}

\begin{exo}[Trois situations]
\explique{On étudie le signe de $d(x)=f(x)-g(x)$ : là où $d(x)>0$, $\Cf$ est au-dessus
de $\mathcal{C}_{g}$.}{méthode 12 du cours}
\cas{a)}
\begin{calculs}
d(x) &= x^{2}-4x+5-(x+1) &&\justif{toute l'expression de $g$}\\
d(x) &= x^{2}-5x+4 &&\\
\Delta &= (-5)^{2}-4\times 1\times 4=9 &&\justif{$\sqrt{9}=3$}\\
x_{1} &= \frac{5-3}{2}=1 &&\\
x_{2} &= \frac{5+3}{2}=4 &&
\end{calculs}
\tabsignes{d(x)}{$-\infty$,$1$,$4$,$+\infty$}{,+,z,-,z,+,}
\explique{$\Cf$ est \res{au-dessous} de $\mathcal{C}_{g}$ sur $\left]1\,;4\right[$,
\res{au-dessus} sur $\left]-\infty\,;1\right[$ et sur $\left]4\,;+\infty\right[$ ; elles se
coupent en $(1\,;2)$ et $(4\,;5)$.}{ordonnées : $g(1)=2$ et $g(4)=5$}
\cas{b)}
\begin{calculs}
d(x) &= -x^{2}+3x+1-(x+1) &&\\
d(x) &= -x^{2}+2x &&\\
d(x) &= -x(x-2) &&\justif{facteur commun $-x$ : racines $0$ et $2$}
\end{calculs}
\tabsignes{d(x)}{$-\infty$,$0$,$2$,$+\infty$}{,-,z,+,z,-,}
\explique{$\Cf$ est \res{au-dessus} de $\mathcal{C}_{g}$ sur $\left]0\,;2\right[$,
\res{au-dessous} ailleurs ; elles se coupent en $(0\,;1)$ et $(2\,;3)$.}{$a=-1<0$ :
positif entre les racines}
\cas{c)}
\begin{calculs}
d(x) &= x^{2}+3x+1-5x &&\\
d(x) &= x^{2}-2x+1 &&\\
\Delta &= (-2)^{2}-4\times 1\times 1=0 &&\justif{une racine double}\\
x_{0} &= \frac{2}{2}=1 &&
\end{calculs}
\tabsignes{d(x)}{$-\infty$,$1$,$+\infty$}{,+,z,+,}
\explique{$\Cf$ est \res{toujours au-dessus} de $\mathcal{C}_{g}$, sauf en $x=1$ où elles
se touchent, au point $(1\,;5)$.}{$d(x)=(x-1)^{2}\ge 0$ ; $g(1)=5$}
\rem{En c), la droite est tangente à la parabole : elle la touche sans la traverser.}
\end{exo}

\begin{exo}[Deux paraboles]
\explique{Même méthode : on étudie le signe de $d(x)=f(x)-g(x)$, en soustrayant
\emph{toute} l'expression de $g$, entre parenthèses.}{méthode 12 du cours}
\cas{a)}
\begin{calculs}
d(x) &= x^{2}+x+1-\left(-x^{2}+4x+3\right) &&\\
d(x) &= x^{2}+x+1+x^{2}-4x-3 &&\justif{le $-$ change tous les signes}\\
d(x) &= 2x^{2}-3x-2 &&\\
\Delta &= (-3)^{2}-4\times 2\times(-2)=25 &&\justif{$9+16$ ; $\sqrt{25}=5$}\\
x_{1} &= \frac{3-5}{4}=-\frac{1}{2} &&\\
x_{2} &= \frac{3+5}{4}=2 &&
\end{calculs}
\tabsignes{d(x)}{$-\infty$,$-\frac{1}{2}$,$2$,$+\infty$}{,+,z,-,z,+,}
\explique{$\Cf$ est \res{au-dessous} de $\mathcal{C}_{g}$ sur
$\left]-\frac{1}{2}\,;2\right[$, \res{au-dessus} ailleurs ; elles se coupent en
$\left(-\frac{1}{2}\,;\frac{3}{4}\right)$ et $(2\,;7)$.}{$a=2>0$ : négatif entre les racines ;
$f\left(-\frac{1}{2}\right)=\frac{1}{4}-\frac{1}{2}+1$ et $f(2)=4+2+1$}
\cas{b)}
\begin{calculs}
d(x) &= 3x^{2}+x-2-\left(3x^{2}-5x+4\right) &&\\
d(x) &= 3x^{2}+x-2-3x^{2}+5x-4 &&\\
d(x) &= 6x-6 &&\justif{les $x^{2}$ s'éliminent}
\end{calculs}
\explique{$6x-6$ s'annule en $1$ ; il est négatif avant, positif après.}{fonction affine
de coefficient $6>0$}
\tabsignes{d(x)}{$-\infty$,$1$,$+\infty$}{,-,z,+,}
\explique{$\Cf$ est \res{au-dessous} de $\mathcal{C}_{g}$ pour $x<1$, \res{au-dessus} pour
$x>1$ ; elles se coupent en $(1\,;2)$.}{$f(1)=3+1-2=2$}
\cas{c)}
\begin{calculs}
d(x) &= 2x^{2}+3x-1-\left(x^{2}+5x+2\right) &&\\
d(x) &= 2x^{2}+3x-1-x^{2}-5x-2 &&\\
d(x) &= x^{2}-2x-3 &&\\
\Delta &= (-2)^{2}-4\times 1\times(-3)=16 &&\justif{$\sqrt{16}=4$}\\
x_{1} &= \frac{2-4}{2}=-1 &&\\
x_{2} &= \frac{2+4}{2}=3 &&
\end{calculs}
\tabsignes{d(x)}{$-\infty$,$-1$,$3$,$+\infty$}{,+,z,-,z,+,}
\explique{$\Cf$ est \res{au-dessous} de $\mathcal{C}_{g}$ sur $\left]-1\,;3\right[$,
\res{au-dessus} ailleurs ; elles se coupent en $(-1\,;-2)$ et $(3\,;26)$.}{$f(-1)=2-3-1$
et $f(3)=18+9-1$}
\rem{Erreur fréquente : $f(x)-g(x)=x^{2}+x+1-x^{2}+4x+3$, en oubliant la parenthèse :
aucun signe n'est changé. La parenthèse autour de $g(x)$ est indispensable.}
\end{exo}

\partie{Problèmes du second degré}

\begin{exo}[Une aire imposée]
\cas{1)}
\explique{Le périmètre vaut $14$~cm, donc longueur $+$ largeur $=7$~cm (le
demi-périmètre). Les deux côtés sont positifs : \res{$0<x<7$}.}{\textbf{inconnue}, et
ses valeurs possibles}
\cas{2)}
\begin{calculs}
\ell &= \res{7-x} &&\justif{largeur : $x+\ell=7$}\\
A(x) &= x(7-x) &&\justif{aire $=$ longueur $\times$ largeur}\\
A(x) &= \res{-x^{2}+7x} &&
\end{calculs}
\cas{3)}
\begin{calculs}
-x^{2}+7x &> 10 &&\justif{l'inéquation}\\
-x^{2}+7x-10 &> 0 &&\justif{tout à gauche}\\
\Delta &= 7^{2}-4\times(-1)\times(-10)=9 &&\justif{$49-40$ ; $\sqrt{9}=3$}\\
x_{1} &= \frac{-7-3}{-2}=5 &&\\
x_{2} &= \frac{-7+3}{-2}=2 &&
\end{calculs}
\tabsignes{-x^{2}+7x-10}{$-\infty$,$2$,$5$,$+\infty$}{,-,z,+,z,-,}
\explique{$a=-1<0$ : positif entre les racines, donc $2<x<5$, qui est bien dans
$\left]0\,;7\right[$.}{\textbf{retour à l'énoncé}}
\explique{L'aire dépasse $10~\text{cm}^{2}$ quand la longueur est \res{strictement
comprise entre $2$~cm et $5$~cm}.}{contrôle : $x=3$ donne $3\times 4=12>10$}
\rem{Si l'on veut que la longueur soit le plus grand côté ($x\ge 7-x$), on ajoute
$x\ge 3{,}5$ : la réponse devient $3{,}5\le x<5$.}
\end{exo}

\begin{exo}[Une trajectoire]
\cas{1)}
\begin{calculs}
h(0) &= -5\times 0^{2}+30\times 0+35=35 &&\justif{au départ, $t=0$}
\end{calculs}
\explique{Le projectile part de \res{$35$~m} de haut.}{la hauteur du toit}
\cas{2)}
\begin{calculs}
\alpha &= -\frac{30}{2\times(-5)}=3 &&\justif{$a=-5<0$ : maximum au sommet}\\
h(3) &= -5\times 9+90+35 &&\\
h(3) &= -45+90+35=80 &&
\end{calculs}
\explique{Hauteur maximale \res{$80$~m}, atteinte au bout de \res{$3$~s}.}{}
\cas{3)}
\begin{calculs}
-5t^{2}+30t+35 &= 0 &&\justif{au sol : $h(t)=0$}\\
t^{2}-6t-7 &= 0 &&\justif{on divise par $-5$}\\
\Delta &= (-6)^{2}-4\times 1\times(-7)=64 &&\justif{$36+28$ ; $\sqrt{64}=8$}\\
t_{1} &= \frac{6-8}{2}=-1 &&\\
t_{2} &= \frac{6+8}{2}=7 &&
\end{calculs}
\explique{Le projectile touche le sol au bout de \res{$7$~s}.}{on rejette $-1$ : le temps
est positif}
\cas{4)}
\begin{calculs}
-5t^{2}+30t+35 &> 35 &&\justif{plus haut que le départ}\\
-5t^{2}+30t &> 0 &&\justif{on retranche $35$}\\
t^{2}-6t &< 0 &&\justif{on divise par $-5<0$ : le sens change}\\
t(t-6) &< 0 &&\justif{racines $0$ et $6$}
\end{calculs}
\explique{$a=1>0$ : négatif entre les racines, donc $t\in\left]0\,;6\right[$. Le
projectile est plus haut que son point de départ pendant \res{$6$~s}.}{$6-0=6$}
\rem{Diviser par un nombre négatif change le sens d'une inégalité, pas celui d'une
égalité.}
\end{exo}

\begin{exo}[Un atelier]
\cas{1)}
\begin{calculs}
B(x) &= R(x)-C(x) &&\justif{bénéfice $=$ recette $-$ coût}\\
B(x) &= 42x-\left(x^{2}-8x+400\right) &&\justif{toute l'expression de $C$}\\
B(x) &= 42x-x^{2}+8x-400 &&\justif{le $-$ change tous les signes}\\
B(x) &= \res{-x^{2}+50x-400} &&
\end{calculs}
\cas{2)}
\begin{calculs}
\Delta &= 50^{2}-4\times(-1)\times(-400) &&\justif{on cherche les racines de $B$}\\
\Delta &= 2\,500-1\,600=900 &&\justif{$\sqrt{900}=30$}\\
x_{1} &= \frac{-50-30}{-2}=40 &&\\
x_{2} &= \frac{-50+30}{-2}=10 &&
\end{calculs}
\tabsignes{B(x)}{$0$,$10$,$40$,$50$}{,-,z,+,z,-,}
\explique{L'atelier fait un bénéfice ($B(x)>0$) pour une production \res{strictement
comprise entre $10$ et $40$ pièces} par jour.}{$a=-1<0$ : positif entre les racines ;
tableau sur $\intff{0}{50}$}
\cas{3)}
\begin{calculs}
\alpha &= -\frac{50}{2\times(-1)}=25 &&\justif{$a<0$ : maximum au sommet}\\
B(25) &= -625+1\,250-400=225 &&\justif{$25^{2}=625$}
\end{calculs}
\explique{Le bénéfice est maximal pour \res{$25$ pièces} par jour ; il vaut alors
\res{$225$~€}.}{$25$ est bien entre $0$ et $50$}
\rem{Contrôle : $25$ est au milieu des racines $10$ et $40$.}
\end{exo}

\end{document}
